\documentclass[11pt]{amsart}
\usepackage{xcolor}
\usepackage[margin=1in]{geometry}
\usepackage{amsmath,amsthm,mathtools}
\usepackage{unicode-math}
\setmathfont{Latin Modern Math}
\usepackage[shortlabels]{enumitem}
\usepackage{setspace}
\usepackage{hyperref}
\usepackage{todonotes}

\usepackage{fontspec}
\directlua{
  luaotfload.add_fallback("leanfallback", {
    "NotoSansMath-Regular:mode=harf;range={0x1D400-0x1D7FF}",
    "STIXGeneral-Regular:mode=harf;range={0x1D400-0x1D7FF}",
    "FreeSerif:mode=harf;range={0x1D400-0x1D7FF}"
  })
}
\setmonofont[
  RawFeature={fallback=leanfallback}
]{JuliaMono}
\usepackage{minted}
\newmintinline[lean]{lean4}{bgcolor=white}
\newminted[leancode]{lean4}{fontsize=\footnotesize,escapeinside=!!}
\usemintedstyle{tango}


\numberwithin{equation}{section}

\newtheorem{theorem}{Theorem}[section]
\newtheorem{proposition}[theorem]{Proposition}
\newtheorem{lemma}[theorem]{Lemma}
\newtheorem{corollary}[theorem]{Corollary}
\theoremstyle{definition}
\newtheorem{definition}[theorem]{Definition}
\theoremstyle{remark}
\newtheorem{remark}[theorem]{Remark}
\newtheorem{question}[theorem]{Question}

\newcommand{\E}{\mathbb E}
\newcommand{\Prob}{\mathbb P}
\newcommand{\N}{\mathbb N}
\newcommand{\R}{\mathbb R}
\newcommand{\Span}{\operatorname{span}}
\newcommand{\supp}{\operatorname{supp}}
\newcommand{\sgn}{\operatorname{sign}}
\newcommand{\conc}{\operatorname{Q}}
\newcommand{\eps}{\varepsilon}
\newcommand{\abs}[1]{\lvert #1 \rvert}
\newcommand{\norm}[1]{\lVert #1 \rVert}
\newcommand{\ip}[2]{\langle #1,#2\rangle}

\title[Stable Phase Retrieval for Independent Coordinates]{Stable Phase Retrieval for Spans of Independent Random Variables}

\author{Pedro Abdalla\textsuperscript{1}}
\author{Jaume de Dios Pont\textsuperscript{2}}
\author{Jo\~ao P. G. Ramos\textsuperscript{3}}
\author{Mitchell A. Taylor\textsuperscript{4}}

\begin{document}
\setstretch{1.2}

\begin{abstract} We prove that, after $L^2$ normalization, stable phase retrieval
holds over the $L^2$-spans of independent real-valued random variables if and only if all but possibly one coordinate satisfies a uniform two-sided
$L^1$ bound. This provides a complete characterization of stable phase retrieval for such spaces, building upon the pioneering work by Calderbank--Daubechies--Freeman--Freeman and answering the main question posed to us by these authors.

We provide two different proofs of this fact. {\color{red} REWRITE: The first is a compactness proof
based on a probe-function-driven profile decomposition, coupled with the inherent infinite divisibility of limit laws of diffuse sums. The second is a quantitative proof, which uses the head-tail decomposition as a template as in the first proof, but replaces the compactness step by an explicit dichotomy through anticoncentration estimates à la Sperner for diffuse tails.} This latter proof was LLM generated based on the ideas in the first proof and important human-based guidance which steered the LLM into producing a quantitative result. The main ingredients in the quantitative proof have been autoformalized in Lean 4. 
\end{abstract}

\maketitle
\begingroup
\renewcommand{\thefootnote}{\arabic{footnote}}
\footnotetext[1]{Department of Mathematics, University of California, Irvine, Irvine, California 92697, USA}
\footnotetext[2]{Center for Data Science, New York University, New York, New York 10011, USA}
\footnotetext[3]{Instituto de Matem\'atica Pura e Aplicada (IMPA), Rio de Janeiro, RJ, Brazil}
\footnotetext[4]{Department of Mathematics, ETH Z\"urich, Z\"urich, Switzerland}
\endgroup
\setcounter{footnote}{0}

\section{Introduction}

\subsection{Historical background} Phase retrieval is a class of inverse problems in which the measurement device captures only magnitude information, while the phase remains unobserved. It dates at least as far back as the seminal Pauli problem in quantum mechanics, which asks how much information about a wave function is encoded in the magnitudes of the function and its Fourier transform \cite{MR3552875,MR3256781,pauli1933allgemeinen,ramos-sousa-pauli}. It also arises naturally in optics and signal processing, where detectors can typically only record  the  intensity of the signal \cite{MR3069958,eldar2014stability,grohs2020survey,sanz1984phase}. In a nutshell, phase retrieval aims to recover a function $f$ while only observing its magnitude $|f|$, up to inherent and acceptable ambiguities. For example, in the real-valued setting, the unavoidable
ambiguity is multiplication by a global sign, and the natural question is whether
this is the only ambiguity.  

The finite-dimensional theory of phase retrieval has been extensively developed in the literature, both for random and structured measurements. The fundamental questions of developing injectivity and stability criteria are by now somewhat well-understood, and practical tools such as polynomial-time algorithms to solve the phase retrieval problem
 have been developed \cite{bandeira2014saving,MR3260258,MR3069958,DB,eldar2014stability,MR3746047,KS}.
In contrast, the infinite-dimensional theory is more delicate and less understood. In fact, injectivity may hold in many Hilbert spaces of interest \cite{cahill2016infinite,freeman2025cahill}, but stability does not follow directly from it. Even worse, uniform stability cannot hold for continuous
frames over infinite dimensional spaces \cite{AG-continuous-frames}, so canonical structured systems
such as the Gabor transform induce ill-posed phase retrieval problems
\cite{alaifari2021gabor,alaifari2025cheeger,grohs2021stable,steinerberger2020stability}. For general background, we refer the reader to \cite{grohs2020survey}.
%Further landmarks include early reconstruction and finite-dimensional stability
%results
%\cite{balan2013,bandeira2014saving,MR3260258,MR3069958,DB,eldar2014stability,MR3746047,KS},
%variants for affine, vector-valued, projection-based models or related nonlinear recovery problems
%\cite{abdalla2025sharp,bartusel2021phase,casazza2017weak,casazza2013,akrami2021note,casazza2017norm,chen2019phase},
%and infinite-dimensional or structured-measurement results
%\cite{localphase,calderbank2022stable,MR4298599}.
%For the surrounding frame-theoretic, function space and discretization background we also refer
%to
%\cite{CHL,DPSTT,FG-atomic,fg22,MR3901674,kashin2021sampling,luef2021gaussian,NOU,MR0423039}.

Mathematically, the injectivity problem asks whether the map $f\mapsto |f|$ is injective on the signal space, after quotienting by the natural global phase ambiguity. The second fundamental question is about stability: does the inverse of this map satisfy a Lipschitz/H\"older stability estimate? Although both questions are relevant for applications, stability is indispensable for any feasible reconstruction algorithm, as injectivity alone is
too weak for numerical reconstruction.

One of the main conceptual contributions of 
\cite{calderbank2022stable,camunez2025characterization,christ2022examples,FOTP,garcia2025isometric,garcia2025existence} is to make the stability question intrinsic to the geometry of the
subspace rather than of a particular discretized frame.  For this reason, we adopt the function-space viewpoint introduced in
\cite{calderbank2022stable,FOTP}.  Let $(X,\|\cdot\|_X)$ be a real normed function space and let
$E\subset X$ be a subspace.  We say that $E$ does \emph{stable phase
retrieval} if there exists $C<\infty$ such that
\begin{equation}
\label{def:stable_pr}
\min_{\sigma\in\{\pm1\}}\norm{f-\sigma g}_X
\le C\,\norm{\abs{f}-\abs{g}}_X
\qquad\text{for all }f,g\in E.
\end{equation}
 
\subsection{Main result} This paper is concerned with the quantitative aspects of stable phase retrieval, by studying random subspaces $E$ spanned
by independent random variables in $L^2$. Such subspaces provide a natural probabilistic analogue of coordinate subspaces and are well suited to understand stable phase retrieval by exploiting the geometry of sums of independent random variables.

In
this setting, Calderbank, Daubechies, Freeman, and Freeman
\cite{calderbank2022stable} proved stable phase retrieval for a simplified model with shifted functions, with the latter two authors proposing to us the following conjecture regarding the characterization of independent sequences doing stable phase retrieval: if $(\xi_i)$ is an orthonormal sequence of independent centered
random variables in $L^2$, then $\overline{\Span}(\xi_i)$ does stable phase
retrieval if and only if there exist constants $0<A\le B<1$ such that
$A\le \norm{\xi_i}_{L^1}\le B$ for all but at most one index. 

The necessity of the two-sided uniform estimates on the $L^1$ norm of $\xi_i$ is known in the literature. In fact, it is shown in \cite{FOTP} that $A$ must be bounded away from zero. Note that this prevents the possibility to construct an almost disjoint sequence in $\overline{\Span}(\xi_i)$, which is a common bottleneck for phase retrieval. The requirement $B<1$ excludes the Rademacher
random variables (random signs) which satisfy $|\xi_i|=1$ almost surely, and therefore violate phase retrieval.

Moreover, the restriction to “at most one exceptional coordinate” is sharp. In fact, two independent Rademacher random variables in the spam are enough to prevent stable phase retrieval to hold for any finite $C$. 

The core question is whether these conditions are also sufficient for stable retrieval to hold. Christ, Pineau and the last-named author \cite{christ2022examples} proved that they suffice under an additional assumption of finite fourth moment. 

The main theorem of this manuscript gives the desired sufficiency result, answering completely the central conjecture in this line and providing an explicit and optimal characterization of when a subspace spanned by independent, centered, real-valued $L^2$-random variables satisfies stable phase retrieval. 

\begin{theorem}[Stable phase retrieval for independent random variables]\label{thm:main}
Let $\eta,\xi_2,\xi_3,\dots$ be independent real-valued random variables and
fix $\delta>0$.  Assume that
\[
\norm{\eta}_{L^2}=1,
\qquad
\E[\eta]=0,
\]
and, for every $i\ge2$,
\[
\norm{\xi_i}_{L^2}=1,
\qquad
\E[\xi_i]=0,
\qquad
\delta\le \norm{\xi_i}_{L^1}\le 1-\delta.
\]
Then there exists $C_{\delta,\eta}\ge1$ such that, for all
$X,Y\in \overline{\Span}(\eta,\xi_i:i\ge2)$,
\[
\min_{\sigma\in\{\pm1\}}\norm{X-\sigma Y}_{L^2}
\le
C_{\delta,\eta}\,\norm{\abs{X}-\abs{Y}}_{L^2}.
\]
\end{theorem}

We present two proofs of this result. The first is a qualitative and conceptually enlightening proof, while the second yields quantitative estimates for $C_{\delta,\eta}$. A common starting point for both proofs is 
a theorem of Freeman, Oikhberg, Pineau, and the last-named author \cite{FOTP}, which establishes the following: for subspaces of $L^2$, it suffices to prove stable phase retrieval for orthogonal pairs (see Proposition \ref{prop:orth-reduction} for the formal statement).

Next, we highlight the main ideas behind our proofs. In the first proof we assume, for the sake of a contradiction, that the result does not hold.  We consider a head-tail decomposition of the coefficients of counterexample-generating pairs of random variables $(\{X_n\}_n, \{Y_n\}_n)$ in the span of $\{\eta,\xi_2,\xi_3,\dots\},$ by splitting them into two components. After rearranging the coordinates, the first component consists of the coordinates that converge strongly in the $\ell_2$ sense. We refer to it as the `head' component. The remaining coordinates may be chosen to be diffuse, i.e., their $\ell^{\infty}$-norms vanish -- hence constituting the `tail' component. 

To bypass the obstacle that the diffuse coordinates may not converge strongly, we relied on characterizations of weak limit of triangular arrays. By independence between the coordinates, we may condition on the head component without affecting the law of the tail component. More accurately, we  derived an equality in the limit for the pair $(\{X_n\}_n, \{Y_n\}_n)$ of the form $|H_X + R_X| = |H_Y + R_Y|,$ where $H_X = \tau H_Y$ for some universal sign $\tau \in \{\pm1\},$ and $(R_X,R_Y)$ is an infinitely divisible random vector that is  independent of $(H_X,H_Y)$. From this equality, it is relatively straightforward to reach a contradiction by conditioning on $(H_X,H_Y)$ and invoking geometric properties of infinite divisible laws.

Since the second proof focus on quantitative estimates for the constant $C_{\delta,\eta}$ in Theorem \ref{thm:main}, we had to modify the argument involving infinite divisible laws as they only appears in the limit as $n$ diverges. To this end, the second proof keeps the same head--tail decomposition of the first proof, but it replaces the limit law of triangular arrays  by an
explicit concentration-diffusion dichotomy for the tail part. A key extra ingredient is an anticoncentration estimate, which acts as a quantitative version of the fact that infinitely divisible laws cannot concentrate on two different lines -- a crucial fact exploited in the first proof.

In that way, the compactness and quantitative proofs address
the same theorem from complementary angles: the compactness proof shows that a
bad sequence cannot exist by passing to a limit model, while the quantitative proof
shows directly that every orthogonal pair already contains a definite amount of
modulus separation.


\subsection{Roadmap}
The rest of this manuscript is organized as follows.   Section~\ref{sec:preliminaries} collects
the basic definitions and tools used throughout this manuscript. 
Section~\ref{sec:compactness-proof}  presents the compactness proof of
Theorem~\ref{thm:main}, including a Gaussian alternative to the final
line-support argument.  Section~\ref{sec:quant-proof} gives a second proof of
Theorem~\ref{thm:main}, explaining how the limit-law step can be bypassed by an
explicit decomposition together with an anticoncentration argument.
Finally, Section~\ref{sec:comments} discusses extensions, related directions,
and open problems.   

\section*{Usage of large language models} 

Large language models played a significant role in the development of this work. Our starting point was a fully human proof of an $L^p$ phase-retrieval result for $p> 2$, based on compactness and central-limit ideas. Our goal, however, was the endpoint $L^2$ statement. By adapting arguments found by GPT in the literature, we managed to find another human-made proof of the endpoint $L^2$ result, by leveraging a non-trivial Gaussian part of the limiting infinitely divisible distributions (see Section \ref{ssec:gaussian-alt} for more details). 

A first key contribution of an LLM was the suggestion of the general principle behind this mechanism: one may forego the Gaussian nature of the limit and replace it by the fact that a nontrivial infinitely divisible random vector in $\R^2$ supported on the cross 
$$\{ (x,y) \in \R^2 \colon |x| = |y|\}$$
must, in fact, be supported on one of the two diagonals $\{x=y\}\cup \{x=-y\}.$ This observation became one of the conceptual pillars of the project. 

A second key contribution came at the quantitative level. Seeking an argument that was both explicitly quantitative and better suited both to uniform estimates as well as to formal verification, we prompted LLMs for an alternative to the infinitely divisible proof. The resulting LLM-generated strategy was combinatorial, of Sperner/antichain type, and yielded a quantitative version of the main result, showing that under the appropriate divisibility hypotheses, closeness to the cross forces closeness to one of the two diagonals. 

This quantitative argument was then formalized and machine-checked in Lean. In this project, GPT was used mainly for mathematical exploration, and Claude Opus 4.6 for assistance with Lean formalization. The authors independently checked all of the final statements and  substantially rewrote the proofs. We take full responsibility for the contents of the paper. 

\section*{Acknowledgments} 

P.A.~was supported by the NSF and by the Simons Collaboration on Theoretical Foundations of Deep Learning. J.P.G.R.~was supported by the FCT through project SHADE (project 2023.17881.ICDT, DOI  10.54499 / 2023.17881.ICDT) and by FAPERJ through the JCNE grant no.~SEI-260003/020475/2025.

\section{Preliminaries}\label{sec:preliminaries}

\subsection{Stable phase retrieval and the orthogonal reduction} We start by recalling the definition of stable phase retrieval in  Equation \ref{def:stable_pr}. Next, we state an important reduction from \cite[Theorem~3.9]{FOTP} or  \cite[Theorem~2.1]{localphase} in the simplified setting of $L^2$ spaces which is our focus here. 

%\begin{definition}
%Let $E$ be a subspace of a real Banach lattice $X$.  We say that $E$ does
%\emph{stable phase retrieval} if there exists $C<\infty$ such that
%\[
%\min_{\sigma\in\{\pm1\}}\norm{f-\sigma g}_X
%\le
%C\,\norm{\abs{f}-\abs{g}}_X
%\qquad(f,g\in E).
%\]
%\end{definition}


\begin{proposition}[Orthogonal reduction]\label{prop:orth-reduction}
Let $E\subset L^2(\mu;\mathbb{R})$ be a subspace. Then $E$
does stable phase retrieval if and only if there exists $c>0$ such that
\[
\norm{\abs{f}-\abs{g}}_{L^2}\ge c
\]
for every orthogonal pair $f,g\in E$ with $\norm{f}_{L^2}=\norm{g}_{L^2}=1$. 
\end{proposition}

\begin{proof} The result from either \cite[Theorem~3.9]{FOTP} or \cite[Theorem~2.1]{localphase} shows that, in order to verify stable phase retrieval, it suffices to only verify it for orthogonal functions $f$ and $g$ such that $\|f\|_{L^2}=1$ and $\|g\|_{L^2}\le 1$. Suppose then that the assertion of Proposition \ref{prop:orth-reduction} holds, but stable phase retrieval does \emph{not} hold. Thus, there exists $\{(f_n,g_n)\}_{n \in \N}$ a sequence of pairs satisfying $\|f_n\|_{L^2} = 1 \geq \|g_n\|_{L^2}$ and $f_n \perp g_n$ such that
\[
\limsup_{n \to \infty} \| |f_n| - |g_n| \|_{L^2} = 0.
\]
We first claim that $\|g_n\|_{L^2} \to 1.$ Indeed, if not, then passing to a subsequence we would be able to estimate 
\[
\||f_n| - |g_n|\|_{L^2} \geq \|f_n\|_{L^2} - \|g_n\|_{L^2} \ge c,
\]
which contradicts the fact that the pair $(f_n,g_n)$ violates stable phase retrieval at infinity. Next,  
\begin{equation}\label{eq:comparison}
\||f_n| - |g_n|\|_{L^2} \geq \| |f_n| - \lambda_n |g_n|\|_{L^2} - \frac{|\lambda_n - 1|}{\lambda_n},
\end{equation}
where $\lambda_n := \frac{1}{\|g_n\|_{L^2}}.$ By hypothesis, the first summand on the right-hand side of \eqref{eq:comparison} is bounded from below by $c$, while the term $\frac{|\lambda_n - 1|}{\lambda_n}$ vanishes because $\|g_n\|_{L^2}$ converges to one. Hence, we reach again a contradiction, which shows that the reduction we were after was indeed valid. 
\end{proof}



\subsection{Probe functions and lower $L^1$ bounds}

The first input is the existence of bounded dual probes.  These are the basic
coefficient-extraction devices in our proofs, as it will become gradually clear in what follows.

\begin{lemma}[Linear and quadratic probes]\label{lem:probes}
Let $\xi$ be a real random variable with $\E[\xi]=0$ and $\E[\xi^2]=1$.

\begin{enumerate}[label=\rm(\arabic*)]
\item If $\E[\abs{\xi}]\ge A$ for some constant $A>0$, then there exists a bounded measurable
function $S$ such that
\[
\E[S]=0,
\qquad
\E[\xi S]=1,
\qquad
\norm{S}_{L^\infty}\le \frac{2}{A}.
\]
\item If $\E[\abs{\xi}]\le B$ for some constant $B<1$, then there exists a bounded measurable
function $T$ such that
\[
\E[T]=0,
\qquad
\E[(\xi^2-1)T]=1,
\qquad
\norm{T}_{L^\infty}\le \frac{2}{1-B}.
\]
\end{enumerate}
\end{lemma}

\begin{proof}
We first construct the linear probe.  Define
\[
S_0:=\sgn(\xi),
\]
with the convention $\sgn(0)=0$.  Then $S_0$ is measurable and satisfies
$|S_0|\le 1$ almost surely.  Moreover,
\[
\xi S_0=\abs{\xi}
\qquad\text{almost surely,}
\]
so
\[
\E[\xi S_0]=\E[\abs{\xi}]\ge A>0.
\]
This positivity allows us to normalize.  Set
\[
S:=\frac{S_0-\E[S_0]}{\E[\xi S_0]}.
\]
The function $S$ is measurable and bounded. In addition,
\[
\E[S]=\frac{\E[S_0]-\E[S_0]}{\E[\xi S_0]}=0,
\]
and, using that $\E[\xi]=0$,
\[
\E[\xi S]
=\frac{\E[\xi S_0]-\E[S_0]\E[\xi]}{\E[\xi S_0]}
=\frac{\E[\xi S_0]}{\E[\xi S_0]}
=1.
\]
Finally, since $|S_0|\le1$ and $|\E[S_0]|\le \E|S_0|\le1$,
\[
|S_0-\E[S_0]|\le 2
\qquad\text{almost surely,}
\]
and therefore
\[
\norm{S}_{L^\infty}
\le \frac{2}{\E[\xi S_0]}
=\frac{2}{\E[\abs{\xi}]}
\le \frac{2}{A}.
\]
This proves part~\rm(1).

For the quadratic probe, apply the same idea to the centered variable
\[
\eta:=\xi^2-1.
\]
Because $\E[\xi^2]=1$, one has $\E[\eta]=0$.  We now show that $\eta$ has
positive $L^1$ mass bounded away from zero.  Since $\xi^2=\abs{\xi}^2$,
\[
\abs{\eta}
=\abs{\abs{\xi}^2-1}
=\abs{\abs{\xi}-1}(\abs{\xi}+1)
\ge \abs{\abs{\xi}-1}.
\]
Taking expectations and using the reverse triangle inequality yields
\[
\E[\abs{\eta}]
\ge \E[\abs{\abs{\xi}-1}]
\ge \bigl|\E[\abs{\xi}]-1\bigr|
=1-\E[\abs{\xi}]
\ge 1-B.
\]
Now set
\[
T_0:=\sgn(\eta)
\]
and define
\[
T:=\frac{T_0-\E[T_0]}{\E[\eta T_0]}
=\frac{T_0-\E[T_0]}{\E[\abs{\eta}]}.
\]
Exactly as in the linear case, $T$ is measurable, bounded, and satisfies
\[
\E[T]=0,
\qquad
\E[\eta T]=1.
\]
Since $\eta=\xi^2-1$, we have
\[
\E[(\xi^2-1)T]=1.
\]
The $L^\infty$ bound is
\[
\norm{T}_{L^\infty}
\le \frac{2}{\E[\abs{\eta}]}
\le \frac{2}{1-B}.
\]
This completes the proof.
\end{proof}

\begin{remark}\label{rem:probe-extraction}
If $\xi_1,\dots,\xi_n$ are independent centered random variables and
$X_c=\sum_i c_i\xi_i$, then the centering of the probes eliminates every
coordinate except the one being tested.  In particular,
\[
\E[X_c\,S_i]=c_i,
\qquad
\E[X_c^2\,T_i]=c_i^2,
\qquad
\E[X_c^2\,S_iS_j]=2c_ic_j
\quad(i\ne j),
\]
whenever $S_i,S_j,T_i$ are the corresponding probes for $\xi_i,\xi_j$.
\end{remark}

The second basic estimate allow us to transfer the lower bounds for the $L^1$ norm of the random variables to its entire spam (up to a different constant). More accurately, we prove the following result.

\begin{lemma}[$L^1$ lower bound for independent sums]\label{lem:l1lower}
Let $\zeta_1,\dots,\zeta_N$ be independent centered real random variables and
assume that
\[
\norm{\zeta_i}_{L^1}\ge \delta \norm{\zeta_i}_{L^2}
\qquad(1\le i\le N).
\]
Then for every scalar family $(c_i)_{i=1}^N$,
\[
\norm{\sum_{i=1}^N c_i\zeta_i}_{L^1}
\ge
\frac{\delta}{2\sqrt2}
\Bigl(\sum_{i=1}^N c_i^2\norm{\zeta_i}_{L^2}^2\Bigr)^{1/2}.
\]
\end{lemma}

\begin{proof}
The statement is homogeneous in the coefficient vector, so we may first
normalize and assume that
\[
\sum_{i=1}^N c_i^2\norm{\zeta_i}_{L^2}^2=1.
\]
Indeed, once the claim is established under this normalization, the general
case follows by scaling.

We next reduce to the unit-variance case.  For each $i$ set
\[
\widetilde\zeta_i:=\frac{\zeta_i}{\norm{\zeta_i}_{L^2}},
\qquad
\widetilde c_i:=c_i\norm{\zeta_i}_{L^2}.
\]
The variables $\widetilde\zeta_i$ are still independent and centered,
they satisfy $\norm{\widetilde\zeta_i}_{L^2}=1$, and the lower $L^1$ bound
becomes
\[
\norm{\widetilde\zeta_i}_{L^1}
=\frac{\norm{\zeta_i}_{L^1}}{\norm{\zeta_i}_{L^2}}
\ge \delta.
\]
Moreover,
\[
\sum_{i=1}^N \widetilde c_i \widetilde\zeta_i
=\sum_{i=1}^N c_i\zeta_i,
\qquad
\sum_{i=1}^N \widetilde c_i^2
=1.
\]
Thus it is enough to prove the lemma under the additional assumptions
\[
\norm{\zeta_i}_{L^2}=1
\qquad(1\le i\le N),
\qquad
\sum_{i=1}^N c_i^2=1.
\]
Define
\[
S:=\sum_{i=1}^N c_i\zeta_i,
\]
and let
\[
S':=\sum_{i=1}^N c_i\zeta_i'
\]
be an independent copy of $S$, obtained from an independent copy
$(\zeta_i')_{i=1}^N$ of the family $(\zeta_i)_{i=1}^N$.
By the triangle inequality,
\[
\abs{S-S'}\le \abs{S}+\abs{S'}.
\]
Taking expectations and using that $S$ and $S'$ have the same distribution, we
obtain
\[
\norm{S-S'}_{L^1}\le 2\norm{S}_{L^1}.
\]
Equivalently,
\[
2\norm{S}_{L^1}\ge \norm{S-S'}_{L^1}.
\]

Now the random variables $\zeta_i-\zeta_i'$ are independent and symmetric.
Conditioning on their absolute values and applying the Khintchine inequality
gives
\[
\norm{S-S'}_{L^1}
=
\Bigl\|\sum_{i=1}^N c_i(\zeta_i-\zeta_i')\Bigr\|_{L^1}
\ge
\frac{1}{\sqrt2}\,
\E\Bigl[\Bigl(\sum_{i=1}^N c_i^2(\zeta_i-\zeta_i')^2\Bigr)^{1/2}\Bigr].
\]

Because the nonnegative numbers $(c_i^2)_{i=1}^N$ sum to one, Cauchy--Schwarz
implies the pointwise estimate
\[
\sum_{i=1}^N c_i^2\abs{\zeta_i-\zeta_i'}
\le
\Bigl(\sum_{i=1}^N c_i^2\Bigr)^{1/2}
\Bigl(\sum_{i=1}^N c_i^2(\zeta_i-\zeta_i')^2\Bigr)^{1/2}
=
\Bigl(\sum_{i=1}^N c_i^2(\zeta_i-\zeta_i')^2\Bigr)^{1/2}.
\]
Taking expectations and combining with the previous inequality, we arrive at
\[
\norm{S-S'}_{L^1}
\ge
\frac{1}{\sqrt2}\sum_{i=1}^N c_i^2\norm{\zeta_i-\zeta_i'}_{L^1}.
\]
Hence
\[
\norm{S}_{L^1}
\ge
\frac{1}{2\sqrt2}\sum_{i=1}^N c_i^2\norm{\zeta_i-\zeta_i'}_{L^1}.
\]

It remains to bound each $\norm{\zeta_i-\zeta_i'}_{L^1}$ from below.  Since
$\zeta_i'$ is independent of $\zeta_i$ and centered,
\[
\E[\zeta_i-\zeta_i' \mid \zeta_i]
=\zeta_i.
\]
Applying Jensen's inequality to the conditional expectation gives
\[
\norm{\zeta_i-\zeta_i'}_{L^1}
\ge
\norm{\E[\zeta_i-\zeta_i' \mid \zeta_i]}_{L^1}
=\norm{\zeta_i}_{L^1}
\ge \delta.
\]
Substituting this into the previous bound yields
\[
\norm{S}_{L^1}
\ge
\frac{\delta}{2\sqrt2}\sum_{i=1}^N c_i^2
=\frac{\delta}{2\sqrt2}.
\]
Since this is exactly the normalized form of the desired estimate, the proof is
complete.
\end{proof}

\subsection{Diffusive arrays and infinitely divisible limits}

A key component in our proof is the rigidity of the weak limits, the next lemma is a precise statement of this type.

\begin{lemma}[Diffuse tail rigidity]\label{lem:tail-rigidity}
Let $(\xi_{n,i})_{n\ge1,\ i\ge2}$ be a triangular array such that for each
fixed $n$ the family $(\xi_{n,i})_{i\ge2}$ is independent, centered,
$L^2$-normalized, and satisfies
\[
\norm{\xi_{n,i}}_{L^1}\ge \delta
\qquad(i\ge2).
\]
Let $c^{(n)}\in \ell^2(\N)$ satisfy
\[
\sup_n \norm{c^{(n)}}_2<\infty,
\qquad
c_1^{(n)}=0,
\qquad
\max_{i\ge2}\abs{c_i^{(n)}}\to0.
\]
Set
\[
Z_n:=\sum_{i\ge2} c_i^{(n)}\xi_{n,i}.
\]
If $Z_n\to0$ in distribution, then $\norm{c^{(n)}}_2\to0$.
\end{lemma}

\begin{proof}
Because the limiting law is the Dirac mass at $0$, the assumption
$Z_n\to0$ in distribution implies that
\[
Z_n\to0
\qquad\text{in probability.}
\]
To pass from convergence in probability to convergence of first moments, we use
uniform integrability.

For each $n$, the variables $(\xi_{n,i})_{i\ge2}$ are independent, centered,
and $L^2$-normalized.  Therefore they form an orthonormal family in $L^2$, and
the random sum $Z_n$ satisfies
\[
\norm{Z_n}_{L^2}^2
=\E\Bigl[\Bigl(\sum_{i\ge2} c_i^{(n)}\xi_{n,i}\Bigr)^2\Bigr]
=\sum_{i\ge2}(c_i^{(n)})^2\E[\xi_{n,i}^2]
=\sum_{i\ge2}(c_i^{(n)})^2
=\norm{c^{(n)}}_2^2.
\]
Since $\sup_n \norm{c^{(n)}}_2<\infty$, the sequence $(Z_n)_n$ is bounded in
$L^2$.  Any $L^2$-bounded family is uniformly integrable in $L^1$, so
$(|Z_n|)_n$ is uniformly integrable.  Because $Z_n\to0$ in probability, we may
therefore conclude that
\[
\norm{Z_n}_{L^1}\longrightarrow0.
\]
We now apply Lemma~\ref{lem:l1lower} to each row of the array.  The assumptions of
that lemma are satisfied because every $\xi_{n,i}$ is centered,
$L^2$-normalized, and obeys $\norm{\xi_{n,i}}_{L^1}\ge\delta$.  Hence
\[
\norm{Z_n}_{L^1}
=\Bigl\|\sum_{i\ge2} c_i^{(n)}\xi_{n,i}\Bigr\|_{L^1}
\ge
\frac{\delta}{2\sqrt2}
\Bigl(\sum_{i\ge2}(c_i^{(n)})^2\Bigr)^{1/2}
=\frac{\delta}{2\sqrt2}\norm{c^{(n)}}_2.
\]
The left-hand side tends to zero, so the right-hand side must also tend to
zero.  Therefore
\[
\norm{c^{(n)}}_2\longrightarrow0,
\]
which is the desired conclusion.
\end{proof}

In order to state our next result, we recall that a $\R^n$-valued random variable $S$ is \emph{infinitely divisible} if, for each positive integer $n$, there exist i.i.d.~random variables $\eta_{n,1},\eta_{n,2},\dots,\eta_{n,n}$ such that 
\[
S \overset{d}{\sim} \eta_{n,1} + \cdots + \eta_{n,n}. 
\]
Infinite divisible distributions are a fundamental object in probability theory due to the classical of limit theorem of triangular arrays. See for instance, \cite[Chapter~IV]{gnedenko-kolmogorov}.
The next result shows that the weak limit of diffuse random vectors is an infinitely divisible random variable. 

\begin{lemma}[Diffuse limits are infinitely divisible]\label{lem:diffuse-id-limit}
Let $d\ge1$, and let $(\eta_{n,i})_{n\ge1,\ i\ge1}$ be a triangular array of
$\R^d$-valued random variables.  Assume that, for each fixed $n$, the variables
$(\eta_{n,i})_{i\ge1}$ are independent and that the row sum
\[
S_n:=\sum_{i\ge1} \eta_{n,i}
\]
is well defined as a limit in probability of its finite partial sums.  Assume
also that the array is infinitesimal, in the sense that for every
$\eps>0$,
\[
\sup_{i\ge1}\Prob\{\norm{\eta_{n,i}}>\eps\}\longrightarrow0.
\]
If $S_n$ converges in distribution to an $\R^d$-valued random variable $S$,
then the law of $S$ is infinitely divisible.
\end{lemma}

\begin{proof}
The finite-row version is classical: if
$U_{n,1},\dots,U_{n,k_n}$ are independent $\R^d$-valued random variables,
\[
\max_{1\le j\le k_n}\Prob\{\norm{U_{n,j}}>\eps\}\longrightarrow0
\qquad(\eps>0),
\]
and
\[
\sum_{j=1}^{k_n}U_{n,j}\Rightarrow U,
\]
then the law of $U$ is infinitely divisible.  This is the classical
Kolmogorov--Khintchine theorem for null arrays; see
\cite[Theorem~9.3]{sato-levy}.  In the one-dimensional case it is also the
classical theorem of Gnedenko--Kolmogorov
\cite[Chapter~IV, Section~24]{gnedenko-kolmogorov}.

It remains only to reduce the present countable-row formulation to that
finite-row theorem.  Since $S_n$ is the limit in probability of its finite
partial sums, for each $n$ we may choose an integer $k_n$ such that
\[
\Prob\left\{
\norm{S_n-\sum_{i=1}^{k_n}\eta_{n,i}}>\frac1n
\right\}\le \frac1n.
\]
Set
\[
S_n^{(k)}:=\sum_{i=1}^{k_n}\eta_{n,i}.
\]
Then $S_n^{(k)}-S_n\to0$ in probability and  hence, by Slutsky's theorem,
\[
S_n^{(k)} = (S_n^{(k)}-S_n)+S_n\Rightarrow S.
\]
For each fixed $n$, the summands $\eta_{n,1},\dots,\eta_{n,k_n}$ are independent,
and the finite array remains infinitesimal because
\[
\max_{1\le i\le k_n}\Prob\{\norm{\eta_{n,i}}>\eps\}
\le
\sup_{i\ge1}\Prob\{\norm{\eta_{n,i}}>\eps\}
\longrightarrow0.
\]
Applying the finite-row theorem gives that the law of $S$ is infinitely
divisible.
\end{proof}


\section{First Proof: Compactness}\label{sec:compactness-proof}

This section proves Theorem~\ref{thm:main} through the first method mentioned in the introduction, consisting of an argument by contradiction through extracting limit properties of sequences failing SPR. 

\subsection{Contradiction array and profile decomposition} We start with the basic setup that failure to satisfy SPR gives. 

\begin{proposition}[Normalized counterexample array]\label{prop:counterexample}
If the conclusion of Theorem~\ref{thm:main} is false, then there exist
\begin{enumerate}[label=\rm(\roman*)]
\item a triangular array $(\xi_{n,i})_{n\ge1,\ i\ge2}$ such that for each
fixed $n$ the family $(\eta,\xi_{n,i})_{i\ge2}$ is independent and every
$\xi_{n,i}$ satisfies
\[
\norm{\xi_{n,i}}_{L^2}=1,
\qquad
\E[\xi_{n,i}]=0,
\qquad
\delta\le \norm{\xi_{n,i}}_{L^1}\le 1-\delta;
\]
\item coefficient vectors $a^{(n)},b^{(n)}\in \ell^2(\N)$ with
\[
\norm{a^{(n)}}_2=\norm{b^{(n)}}_2=1,
\qquad
\ip{a^{(n)}}{b^{(n)}}=0,
\]
such that, writing
\[
X_n:=a_1^{(n)}\eta+\sum_{i\ge2} a_i^{(n)}\xi_{n,i},
\qquad
Y_n:=b_1^{(n)}\eta+\sum_{i\ge2} b_i^{(n)}\xi_{n,i},
\]
one has
\[
\norm{\abs{X_n}-\abs{Y_n}}_{L^2}\longrightarrow 0.
\]
\end{enumerate}
\end{proposition}

\begin{proof}
Assume that the conclusion of Theorem~\ref{thm:main} is false.  Then there is
no constant $C$ for which the stable phase retrieval estimate holds uniformly
on the closed span of the coordinates.  By the orthogonal reduction
(Proposition~\ref{prop:orth-reduction}), this means that we may find, for each
$n\in\N$, an orthogonal pair of unit vectors in the span whose modulus gap is
at most $1/n$.  Concretely, there exist functions $X_n,Y_n$ in
$\overline{\Span}(\eta,\xi_i:i\ge2)$ such that
\[
\norm{X_n}_{L^2}=\norm{Y_n}_{L^2}=1,
\qquad
\E[X_nY_n]=0,
\qquad
\norm{|X_n|-|Y_n|}_{L^2}\le \frac1n.
\]

\noindent Because the family $(\eta,\xi_i)_{i\ge2}$ is orthonormal in $L^2$, we may write
\[
X_n=a_1^{(n)}\eta+\sum_{i\ge2}a_i^{(n)}\xi_i,
\qquad
Y_n=b_1^{(n)}\eta+\sum_{i\ge2}b_i^{(n)}\xi_i,
\]
with coefficient vectors $a^{(n)},b^{(n)}\in\ell^2(\N)$.  Parseval's identity
then yields
\[
\norm{a^{(n)}}_2=\norm{X_n}_{L^2}=1,
\qquad
\norm{b^{(n)}}_2=\norm{Y_n}_{L^2}=1,
\]
and, since the two functions are orthogonal in $L^2$,
\[
\ip{a^{(n)}}{b^{(n)}}=\E[X_nY_n]=0.
\]

Finally, to prepare for later rowwise permutations of the good coordinates (the ones with $L^1$ norm bounded away from one), we
rename the good variables used in the $n$-th row by writing them as
$(\xi_{n,i})_{i\ge2}$.  This does not alter their laws or any of the stated
hypotheses, but it is convenient for the compactness scheme.  The displayed
modulus convergence is exactly our chosen violating sequence.  This proves the
proposition.
\end{proof}

\begin{lemma}[Reordered profile decomposition]\label{lem:profile}
After passing to a subsequence and, row by row, applying the same permutation
to the good coordinates and to their coefficients, one may find vectors
$a,b\in\ell^2(\N)$ and integers $m_n\to\infty$ such that, with
\[
a_{\mathrm{head}}^{(n)}:=a^{(n)}\mathbf 1_{\{1,\dots,m_n\}},
\qquad
a_{\mathrm{tail}}^{(n)}:=a^{(n)}\mathbf 1_{\{m_n+1,m_n+2,\dots\}},
\]
and likewise for $b^{(n)}$, the following hold:
\begin{enumerate}[label=\rm(\arabic*)]
\item $a_{\mathrm{head}}^{(n)}\to a$ and $b_{\mathrm{head}}^{(n)}\to b$
strongly in $\ell^2$;
\item
\[
\max_{i>m_n}\max(\abs{a_i^{(n)}},\abs{b_i^{(n)}})\to0;
\]
\item
\[
\norm{a_{\mathrm{tail}}^{(n)}}_2^2\to1-\norm{a}_2^2,
\qquad
\norm{b_{\mathrm{tail}}^{(n)}}_2^2\to1-\norm{b}_2^2,
\]
and
\[
\ip{a_{\mathrm{tail}}^{(n)}}{b_{\mathrm{tail}}^{(n)}}
\to -\ip{a}{b}.
\]
\end{enumerate}
\end{lemma}

\begin{proof}
For each row $n$ and each good index $i\ge2$, define
\[
s_i^{(n)}:=\max\bigl(\abs{a_i^{(n)}},\abs{b_i^{(n)}}\bigr).
\]
Permuting the good coordinates within the $n$-th row, we may assume that the
sequence $(s_i^{(n)})_{i\ge2}$ is nonincreasing in $i$.  The same permutation
is applied simultaneously to the variables $\xi_{n,i}$ and to both coefficient
vectors, so our hypothesis is not changed.

The monotone rearrangement gives a uniform pointwise bound on the reordered
coefficients.  Indeed,
\[
\sum_{i\ge2}(s_i^{(n)})^2
\le
\sum_{i\ge2}\bigl((a_i^{(n)})^2+(b_i^{(n)})^2\bigr)
\le 2.
\]
Since $(s_i^{(n)})_{i\ge2}$ is nonincreasing, for every $i\ge2$ we have
\[
(i-1)(s_i^{(n)})^2
\le \sum_{j=2}^{i}(s_j^{(n)})^2
\le 2,
\]
and therefore
\[
s_i^{(n)}\le \sqrt{\frac{2}{i-1}}.
\]
In particular, for each fixed coordinate $i$, the sequences
$\bigl(a_i^{(n)}\bigr)_{n\ge1}$ and $\bigl(b_i^{(n)}\bigr)_{n\ge1}$ are
bounded.

We now extract coordinatewise limits.  First pass to a subsequence so that
$a_1^{(n)}$ and $b_1^{(n)}$ converge.  Next pass to a further subsequence so
that $a_2^{(n)}$ and $b_2^{(n)}$ converge, and continue inductively.  A
standard diagonal argument yields a single subsequence, not relabeled, for
which every fixed coordinate converges:
\[
a_i^{(n)}\to a_i,
\qquad
b_i^{(n)}\to b_i
\qquad(i\ge1).
\]
Set
\[
a:=(a_i)_{i\ge1},
\qquad
b:=(b_i)_{i\ge1}.
\]

We claim that $a,b\in\ell^2(\N)$.  Fix $M\in\N$.  Since the first $M$
coordinates converge,
\[
\sum_{i=1}^{M} a_i^2
=\lim_{n\to\infty}\sum_{i=1}^{M}(a_i^{(n)})^2
\le \liminf_{n\to\infty}\sum_{i\ge1}(a_i^{(n)})^2
=1.
\]
Letting $M\to\infty$ and using monotone convergence gives
$\sum_i a_i^2\le1$.  The same argument yields $\sum_i b_i^2\le1$.

Because $a,b\in\ell^2$, we may choose an increasing sequence of integers
$(m_k)_{k\ge1}$ with $m_k\to\infty$ such that
\[
\sum_{i>m_k} a_i^2+\sum_{i>m_k} b_i^2\le 2^{-k}
\qquad(k\ge1).
\]
For each fixed $k$, coordinatewise convergence implies
\[
\sum_{i=1}^{m_k}\abs{a_i^{(n)}-a_i}^2
+
\sum_{i=1}^{m_k}\abs{b_i^{(n)}-b_i}^2
\longrightarrow0
\qquad(n\to\infty).
\]
Passing to another subsequence if necessary, we may assume that for each $n$,
\[
\sum_{i=1}^{m_n}\abs{a_i^{(n)}-a_i}^2
+
\sum_{i=1}^{m_n}\abs{b_i^{(n)}-b_i}^2
\le 2^{-n}.
\]

We now verify the three conclusions.  First,
\[
\norm{a_{\mathrm{head}}^{(n)}-a}_2^2
\le
\sum_{i=1}^{m_n}\abs{a_i^{(n)}-a_i}^2
+
\sum_{i>m_n}a_i^2
\le 2^{-n}+2^{-n},
\]
so $a_{\mathrm{head}}^{(n)}\to a$ in $\ell^2$.  The same estimate gives
$b_{\mathrm{head}}^{(n)}\to b$.

Second, if $i>m_n$, then by monotonicity,
\[
\max\bigl(\abs{a_i^{(n)}},\abs{b_i^{(n)}}\bigr)
=s_i^{(n)}
\le s_{m_n+1}^{(n)}
\le \sqrt{\frac{2}{m_n}}.
\]
Since $m_n\to\infty$, this proves
\[
\max_{i>m_n}\max\bigl(\abs{a_i^{(n)}},\abs{b_i^{(n)}}\bigr)\to0.
\]

Finally, because each full coefficient vector has $\ell^2$ norm one,
\[
\norm{a_{\mathrm{tail}}^{(n)}}_2^2
=1-\norm{a_{\mathrm{head}}^{(n)}}_2^2
\longrightarrow 1-\norm{a}_2^2,
\]
and similarly
\[
\norm{b_{\mathrm{tail}}^{(n)}}_2^2
\longrightarrow 1-\norm{b}_2^2.
\]
For the inner products, use the orthogonality of the full vectors:
\[
0
=\ip{a^{(n)}}{b^{(n)}}
=\ip{a_{\mathrm{head}}^{(n)}}{b_{\mathrm{head}}^{(n)}}
+\ip{a_{\mathrm{tail}}^{(n)}}{b_{\mathrm{tail}}^{(n)}}.
\]
Since the head component converges strongly,
\[
\ip{a_{\mathrm{head}}^{(n)}}{b_{\mathrm{head}}^{(n)}}
\to \ip{a}{b},
\]
and therefore
\[
\ip{a_{\mathrm{tail}}^{(n)}}{b_{\mathrm{tail}}^{(n)}}
\to -\ip{a}{b}.
\]
This proves the lemma.
\end{proof}

\subsection{Head identification} Our next step is to identify the  behavior of the head profiles constructed in the previous step when $n$ is sufficiently large. To this end, we start with an analogue of Lemma \ref{lem:probes} for the exceptional coordinate $\eta$.

\begin{lemma}[Exceptional probes]\label{prop:exceptional-probes}
There exists a bounded measurable function $S_\eta$ such that
\[
\E[S_\eta]=0,
\qquad
\E[\eta S_\eta]=1.
\]
If $\eta^2\not\equiv1$ almost surely, then there exists a bounded measurable
function $T_\eta$ such that
\[
\E[T_\eta]=0,
\qquad
\E[(\eta^2-1)T_\eta]=1.
\]
\end{lemma}

\begin{proof}
Since $\norm{\eta}_{L^2}=1$, the random variable $\eta$ is not almost surely
zero.  Hence $\E[\abs{\eta}]>0$.  Applying the linear-probe construction from
Lemma~\ref{lem:probes} with $A=\E[\abs{\eta}]$ produces a bounded measurable
function $S_\eta$ satisfying
\[
\E[S_\eta]=0,
\qquad
\E[\eta S_\eta]=1.
\]

Assume now that $\eta^2\not\equiv1$ almost surely.  Then the centered random
variable $\eta^2-1$ is not almost surely zero, because
\[
\E[\eta^2-1]=\E[\eta^2]-1=0.
\]
Therefore its $L^1$ norm is strictly positive:
\[
\E[\abs{\eta^2-1}]>0.
\]
Applying again the linear-probe construction, now to the random variable
$\eta^2-1$, we obtain a bounded measurable function $T_\eta$ such that
\[
\E[T_\eta]=0,
\qquad
\E[(\eta^2-1)T_\eta]=1.
\]
This is exactly the desired conclusion.
\end{proof}

\begin{lemma}[Identification of the head profile]\label{lem:head-sign}
Either there exists $\tau\in\{\pm1\}$ such that
\[
a=\tau b
\qquad\text{and}\qquad
\norm{a_{\mathrm{head}}^{(n)}-\tau b_{\mathrm{head}}^{(n)}}_2\to0,
\]
or else $\eta$ is a Rademacher variable and
\[
a_i=b_i=0
\qquad(i\ge2).
\]
\end{lemma}

\begin{proof}
We first convert the convergence of moduli into convergence of the quadratic
observable. Since
\[
\norm{X_n}_{L^2}=\norm{Y_n}_{L^2}=1,
\]
we have
\[
\norm{|X_n|+|Y_n|}_{L^2}\le 2.
\]
Therefore
\[
\norm{|X_n|^2-|Y_n|^2}_{L^1}
=\norm{(|X_n|-|Y_n|)(|X_n|+|Y_n|)}_{L^1}
\le 2\norm{|X_n|-|Y_n|}_{L^2}
\longrightarrow0.
\]
Because the variables are real-valued, $|X_n|^2=X_n^2$ and $|Y_n|^2=Y_n^2$,
so we have established
\[
\norm{X_n^2-Y_n^2}_{L^1}\longrightarrow0.
\]

\smallskip
\noindent\emph{Step 1: diagonal good-coordinate identities.}
Fix $i\ge2$ and let $T_{n,i}$ be the quadratic probe associated with the
coordinate $\xi_{n,i}$.  The probes are uniformly bounded in $L^\infty$, so the
$L^1$ convergence above implies
\[
\E[(X_n^2-Y_n^2)T_{n,i}]\longrightarrow0.
\]
We compute the left-hand side explicitly.  Expanding $X_n^2$ and integrating
against $T_{n,i}$, every term involving a coordinate different from
$\xi_{n,i}$ vanishes because $T_{n,i}$ depends only on $\xi_{n,i}$ and is
centered.  Mixed terms involving $\xi_{n,i}$ and another coordinate also
vanish by independence and centering.  The only surviving contribution is the
term $(a_i^{(n)})^2\xi_{n,i}^2$.  Hence
\[
\E[X_n^2T_{n,i}]
=(a_i^{(n)})^2\E[\xi_{n,i}^2T_{n,i}]
=(a_i^{(n)})^2,
\]
because
\[
\E[\xi_{n,i}^2T_{n,i}]
=\E[(\xi_{n,i}^2-1)T_{n,i}]+\E[T_{n,i}]
=1+0.
\]
Similarly,
\[
\E[Y_n^2T_{n,i}]=(b_i^{(n)})^2.
\]
Therefore
\[
(a_i^{(n)})^2-(b_i^{(n)})^2
=\E[(X_n^2-Y_n^2)T_{n,i}]
\longrightarrow0.
\]
Passing to the limit in $n$ gives
\[
a_i^2=b_i^2
\qquad(i\ge2).
\]

\smallskip
\noindent\emph{Step 2: off-diagonal good-good identities.}
Fix distinct indices $i,j\ge2$ and let $S_{n,i},S_{n,j}$ be the corresponding
linear probes.  Again boundedness and $L^1$ convergence imply
\[
\E[(X_n^2-Y_n^2)S_{n,i}S_{n,j}]\longrightarrow0.
\]
In the expansion of $X_n^2$, all terms vanish after integration against
$S_{n,i}S_{n,j}$ except the mixed term
$2a_i^{(n)}a_j^{(n)}\xi_{n,i}\xi_{n,j}$.  Indeed, any term missing one of the
coordinates $\xi_{n,i}$ or $\xi_{n,j}$ is killed by the zero mean of the
corresponding probe, and terms involving other coordinates factor through zero
means by independence.  Hence
\[
\E[X_n^2S_{n,i}S_{n,j}]
=2a_i^{(n)}a_j^{(n)}
\E[\xi_{n,i}S_{n,i}]\,\E[\xi_{n,j}S_{n,j}]
=2a_i^{(n)}a_j^{(n)}.
\]
Similarly,
\[
\E[Y_n^2S_{n,i}S_{n,j}]
=2b_i^{(n)}b_j^{(n)}.
\]
Thus
\[
2(a_i^{(n)}a_j^{(n)}-b_i^{(n)}b_j^{(n)})\longrightarrow0,
\]
and passing to the limit gives
\[
a_ia_j=b_ib_j
\qquad(i\ne j,\ i,j\ge2).
\]

\smallskip
\noindent\emph{Step 3: exceptional-good identities.}
Fix $j\ge2$.  Let $S_\eta$ be the bounded probe from
Proposition~\ref{prop:exceptional-probes}.  By the same reasoning,
\[
\E[(X_n^2-Y_n^2)S_\eta S_{n,j}]\longrightarrow0.
\]
Once more only one mixed term survives, namely
$2a_1^{(n)}a_j^{(n)}\eta\xi_{n,j}$ for $X_n^2$ and the corresponding term for
$Y_n^2$.  Therefore
\[
\E[X_n^2S_\eta S_{n,j}]
=2a_1^{(n)}a_j^{(n)}\E[\eta S_\eta]\E[\xi_{n,j}S_{n,j}]
=2a_1^{(n)}a_j^{(n)},
\]
and analogously
\[
\E[Y_n^2S_\eta S_{n,j}]
=2b_1^{(n)}b_j^{(n)}.
\]
Hence
\[
a_1a_j=b_1b_j
\qquad(j\ge2).
\]

\smallskip
\noindent\emph{Step 4: deduction of the sign.}
Assume first that there exists some index $i_0\ge2$ such that $a_{i_0}\ne0$.
From the diagonal identity we know that
$a_{i_0}^2=b_{i_0}^2$, so $b_{i_0}=\tau a_{i_0}$ for some
$\tau\in\{\pm1\}$.  Then for every $j\ge2$, the off-diagonal identity gives
\[
a_{i_0}a_j=b_{i_0}b_j=\tau a_{i_0}b_j.
\]
Since $a_{i_0}\ne0$, we conclude that $a_j=\tau b_j$ for every $j\ge2$.
Applying the exceptional-good identity with $j=i_0$, we obtain
\[
a_1a_{i_0}=b_1b_{i_0}=\tau a_{i_0}b_1,
\]
and again $a_{i_0}\ne0$ implies $a_1=\tau b_1$.  Hence $a=\tau b$.

Assume next that
\[
a_i=0
\qquad\text{for every }i\ge2.
\]
Then the diagonal identities show that also $b_i=0$ for every $i\ge2$.
If $\eta^2\not\equiv1$ almost surely, let $T_\eta$ be the quadratic probe from
Lemma~\ref{prop:exceptional-probes}.  Repeating the diagonal argument for
the exceptional coordinate gives
\[
(a_1^{(n)})^2-(b_1^{(n)})^2
=\E[(X_n^2-Y_n^2)T_\eta]
\longrightarrow0,
\]
whence $a_1^2=b_1^2$ and therefore $a=\tau b$ for some sign $\tau$.

The only remaining possibility is that $\eta^2\equiv1$ almost surely.  Since
$\E[\eta]=0$ and $\norm{\eta}_{L^2}=1$, this means exactly that $\eta$ is a
Rademacher variable.  In that case we are in the second alternative stated in
the lemma.

Finally, if the first alternative holds, then the convergence of the head
profiles follows from Lemma~\ref{lem:profile}:
\[
\norm{a_{\mathrm{head}}^{(n)}-\tau b_{\mathrm{head}}^{(n)}}_2
\le
\norm{a_{\mathrm{head}}^{(n)}-a}_2
+
\norm{b_{\mathrm{head}}^{(n)}-b}_2
\longrightarrow0.
\]
This finishes the proof.
\end{proof}

\subsection{Limit law and exact modulus identity} We next identify the properties that the limiting distributions stemming from sequences failing SPR must satisfy. 

\begin{lemma}[Limit law]\label{lem:limit-law}
After passing to a subsequence, the quadruple
\[
\bigl(X_n^{\mathrm{head}},Y_n^{\mathrm{head}},
X_n^{\mathrm{tail}},Y_n^{\mathrm{tail}}\bigr)
\]
converges in distribution to a random vector
\[
(H_X,H_Y,R_X,R_Y)\in\R^4
\]
such that:
\begin{enumerate}[label=\rm(\arabic*)]
\item $(H_X,H_Y)$ is independent of $(R_X,R_Y)$;
\item the law of $(R_X,R_Y)$ is infinitely divisible;
\item either $H_X=\tau H_Y$ almost surely for some $\tau\in\{\pm1\}$, or
else $\eta$ is Rademacher and
$H_X=a_1\eta$, $H_Y=b_1\eta$ almost surely;
\item
\[
\abs{H_X+R_X}=\abs{H_Y+R_Y}
\qquad\text{almost surely.}
\]
\end{enumerate}
\end{lemma}

\begin{proof}
The head variables converge in $L^2$ by the strong convergence of the head
coefficients and the orthonormality of each row.  The tail coefficient arrays
are infinitesimal by Lemma~\ref{lem:profile}, so every weak limit of the
tail pair is infinitely divisible by Lemma \ref{lem:diffuse-id-limit}.  Independence of head and tail persists
because they are built from disjoint independent coordinates.  Finally,
$\norm{\abs{X_n}-\abs{Y_n}}_{L^2}\to0$ implies
$\abs{X_n}-\abs{Y_n}\to0$ in probability; passing to the joint weak limit
gives the exact identity in the limit.  The description of the head follows
from Lemma~\ref{lem:head-sign}.
\end{proof}

\begin{lemma}[Infinitely divisible laws on the phase-retrieval cone are
line-supported]\label{lem:id-line}
Let $\mu$ be an infinitely divisible probability measure on $\R^2$ whose support satisfies
\[
\supp\mu\subset\{(u,v)\in\R^2:\abs{u}=\abs{v}\} \eqqcolon \Gamma.
\]
Then there exists $\tau\in\{\pm1\}$ such that
\[
\supp\mu\subset \{(u,\tau u):u\in\R\} \eqqcolon L_\tau.
\]
\end{lemma}

\begin{proof}
The set $\Gamma$ is exactly the union of two lines $L_+\cup L_-$.  Since $\mu$ is
infinitely divisible, it is in particular $2$-divisible: there exists a
probability measure $\nu$ such that $\mu=\nu*\nu$.  Let $A=\supp\nu$.  Then
\[
\supp\mu=\overline{A+A}\subset L_+\cup L_-.
\]
If $x,y\in A$ are distinct, then $2x,x+y,2y\in A+A$.  If $2x$ and $2y$ belong
to the same line $L_\tau$, then so does $x+y$, forcing $x,y\in L_\tau$.  If
they belong to different lines, then the midpoint $x+y$ is the origin, so
$y=-x$ and again $x,y$ lie in one of the lines.  Thus every two points of $A$
lie on a common phase-retrieval line.  Since $L_+\cap L_-=\{0\}$, this implies
that $A$ is contained in one of the two lines, and so is $\supp\mu$.
\end{proof}

\subsection{Conclusion of the compactness proof}

\begin{proof}[Proof of Theorem~\ref{thm:main}]
Assume that Theorem~\ref{thm:main} fails and let the counterexample array be as
in Proposition~\ref{prop:counterexample}.  Apply
Lemmas~\ref{lem:profile}, \ref{lem:head-sign}, and \ref{lem:limit-law}.

If the first alternative of Lemma~\ref{lem:head-sign} holds, fix
$\tau\in\{\pm1\}$ such that $H_Y=\tau H_X$ almost surely.  The exact modulus
identity becomes
\[
\abs{H_X+R_X}=\abs{\tau H_X+R_Y}
\qquad\text{almost surely.}
\]
If $H_X$ is not almost surely zero, then conditioning on the tail and varying
the head over two support points of $H_X$ shows that
$R_Y-\tau R_X=0$ almost surely.  If $H_X\equiv0$, then
$(R_X,R_Y)$ itself satisfies $\abs{R_X}=\abs{R_Y}$ almost surely, so
Lemma~\ref{lem:id-line} yields
$R_Y=\tau_0 R_X$ almost surely for some $\tau_0\in\{\pm1\}$.

In either case, there exists a sign $\sigma\in\{\pm1\}$ for which the random variables
\[
Z_n:=\sum_{i>m_n}(b_i^{(n)}-\sigma a_i^{(n)})\xi_{n,i}
\]
converge to $0$ in distribution.  Since the coefficients are diffuse,
Lemma~\ref{lem:tail-rigidity} gives
\[
\norm{b_{\mathrm{tail}}^{(n)}-\sigma a_{\mathrm{tail}}^{(n)}}_2\to0.
\]
Together with Lemma~\ref{lem:head-sign} this implies
$\norm{b^{(n)}-\sigma a^{(n)}}_2\to0$, contradicting
\[
\norm{b^{(n)}-\sigma a^{(n)}}_2^2
=\norm{a^{(n)}}_2^2+\norm{b^{(n)}}_2^2
-2\sigma\ip{a^{(n)}}{b^{(n)}}=2.
\]

If the second alternative of Lemma~\ref{lem:head-sign} holds, then
$\eta$ is Rademacher and the head limits are multiples of $\eta$ alone.  The
same argument applies, since in that case the good coordinates disappear from
the head and the tail limit still lies on $\Gamma$.
\end{proof}

\subsection{A Gaussian alternative in the degenerate-head regime}\label{ssec:gaussian-alt}

The previous endpoint used only the support structure of infinitely divisible
limits.  The discussion below records a second route, showing that in the degenerate-head regime the
diffuse tail automatically generates a nontrivial Gaussian component.

We start with a result that ensures that, in spite of the failure of the Lindeberg condition for the Central Limit Theorem, one may still ensure the existence of a non-degenerated Gaussian component of the limiting distribution. 

\begin{proposition}[Diffuse tails force a nontrivial Gaussian part]
\label{prop:gaussian-part}
Let
\[
V_{n,i}:=(a_i^{(n)}\xi_{n,i},\,b_i^{(n)}\xi_{n,i})\in\R^2
\qquad(i>m_n),
\]
where the coordinates are independent, centered, $L^2$-normalized, and satisfy
$\norm{\xi_{n,i}}_{L^1}\ge \delta$.  Set
\[
\alpha_{n,i}:=\bigl((a_i^{(n)})^2+(b_i^{(n)})^2\bigr)^{1/2}.
\]
Assume
\[
\max_{i>m_n}\alpha_{n,i}\to0,
\qquad
\sum_{i>m_n}\alpha_{n,i}^2\to \sigma^2>0,
\]
and that the row sums $\sum_{i>m_n}V_{n,i}$ converge in distribution to an
infinitely divisible law $\mu$ on $\R^2$.  Then the covariance matrix of the Gaussian component of $\mu$ is not identically zero.
\end{proposition}

\begin{proof}[Sketch of proof]
Set $\beta_n=\max_{i>m_n}\alpha_{n,i}$ and $r_n=\beta_n^{1/2}$.  For each
$i>m_n$ with $\alpha_{n,i}\ne0$, let $K_{n,i}=r_n/\alpha_{n,i}\to\infty$
uniformly and define the centered truncation
\[
\widetilde\xi_{n,i}
:=
\xi_{n,i}\mathbf 1_{\{\abs{\xi_{n,i}}\le K_{n,i}\}}
-\E[\xi_{n,i}\mathbf 1_{\{\abs{\xi_{n,i}}\le K_{n,i}\}}].
\]
Since $\E[\abs{\xi_{n,i}}]\ge\delta$ and $\E[\xi_{n,i}^2]=1$, a simple estimate shows that there exists $c_0=c_0(\delta)>0$ such that
\[
\E[\widetilde\xi_{n,i}^2]\ge c_0
\]
for all large $n$ and all $i>m_n$.  Define
$\widetilde V_{n,i}=(a_i^{(n)}\widetilde\xi_{n,i},\,b_i^{(n)}\widetilde\xi_{n,i})$.
Then $\abs{\widetilde V_{n,i}}\le 2r_n$, so the truncated array is bounded and
infinitesimal, and
\[
\sum_{i>m_n}\E[\abs{\widetilde V_{n,i}}^2]
\ge c_0\sum_{i>m_n}\alpha_{n,i}^2
\longrightarrow c_0\sigma^2>0.
\]
For infinitesimal triangular arrays, the covariance matrices of shrinking
centered truncations converge to the Gaussian covariance of the limit law; see
\cite[Chapter~IV]{gnedenko-kolmogorov}.  Since the traces of the truncation
covariances stay bounded away from zero, the Gaussian covariance of $\mu$ is
nonzero.
\end{proof}

We now show an alternative route to the conclusion of the compactness proof of Theorem \ref{thm:main}. Assume that we are in the degenerate-head regime $H_X\equiv H_Y\equiv0$ -- as the other one is essentially trivial -- so the
tail component $\mu$, the law of $(R_X,R_Y)$, is supported on $\Gamma=\{(u,v):\abs{u}=\abs{v}\}$.
Write $\mu=\gamma*\nu$, where $\gamma$ is the Gaussian part and $\nu$ the
independent remainder term.  By Proposition~\ref{prop:gaussian-part}, $\gamma$ is
nontrivial.  Its support is a nonzero linear subspace $E\subset\R^2$.

Because $\supp\mu=\supp\gamma+\supp\nu\subset\Gamma=L_+\cup L_-$, the
connected set $E$ must itself lie in one of the lines $L_\tau$.  Fix such a
$\tau$.  For any $z\in\supp\nu$, the translated line $E+z$ is contained in
$\Gamma$.  Since a translate of $L_\tau$ not equal to $L_\tau$ intersects
$L_+\cup L_-$ in at most two points, one must have $z\in L_\tau$.  Hence
$\supp\nu\subset L_\tau$ and therefore $\supp\mu\subset L_\tau$.  Thus
$R_Y=\tau R_X$ almost surely, and hence this recovers the final contradiction in the first proof of Theorem \ref{thm:main}. 

\section{Second Proof: Anticoncentration}\label{sec:quant-proof}

This section gives a second proof of Theorem~\ref{thm:main}.  The argument is
quantitative in spirit.  Instead of passing to a limit law, we split the
coefficients into a head and a tail and show directly that one of two explicit
mechanisms forces a positive modulus gap.

By Proposition~\ref{prop:orth-reduction}, it is enough to prove an orthogonal
lower bound for arbitrary unit vectors in the one-exception span. By continuity, it suffices to consider the case with only finitely ($n$) many random variables, as long as the constants do not depend on $n$.
\todo[inline]{Wote an exact theorem that we can cite below, this is what Lean proves, but we can adapt the Lean. I have also simplifed the lower bound significantly}
\begin{proposition}
    Let $\xi_0,\dots, \xi_{n-1}$ be independent random variables, satisfying
    \[
    \forall i = 0, \dots {n-1},\quad \|\xi_i\|_2 = 1, \quad \mathbb E[\xi_i] = 0 \quad  \text{and} \quad \E[|\xi_i|] \ge A>0 
    \]
    \[
    \forall i = 1,\dots n-1, \quad \E[|\xi_i|] \le B<1 
    \]
    Let
    $a=(a_i)_{i\ge1}, b=(b_i)_{i\ge1}\in \ell^2(\N)$ satisfy
    \[
    \norm{a}_2=\norm{b}_2=1,
    \qquad
    \ip{a}{b}=0,
    \]
    and set
    \[
    X:=a_1\eta+\sum_{i\ge2} a_i\xi_i,
    \qquad
    Y:=b_1\eta+\sum_{i\ge2} b_i\xi_i.
    \]
    Then 
    \[
    \left\||\sum_{i\ge2} a_i\xi_i|-|\sum_{i\ge2} b_i\xi_i|\right \|_{L^2}\ge c_{A,B}>0,
    \]
    where for some constant $c_{A,B}$ satisfying
    \[
    c_{A,B}
    \ge 
    2^{-57}
    \min \left(A^{10},A^{8} (1-B)\right)
    \]
\end{proposition}

\todo[inline]{The dependency in $A$ is a disaster but the dependency in $B$ may be sharp? 

I do not think so, a sparse Rademacher gives $\sqrt{1-B}$ which shows that $1-B$ cannot be sharp}



\subsection{The dichotomy setup}

We start with the basic setup for our second proof that is based on a simple dichotomy. Let
\[
\eps:=\E[\abs{X^2-Y^2}],
\qquad
\Delta:=\norm{|X|-|Y|}_{L^2},
\]
and notice that
\[
\eps\le 2\Delta.
\]
Fix
\[
\kappa_{\delta,\eta}
:=
\min\left(
\frac{\delta}{4\sqrt2},
\frac{\E[\abs{\eta-\widetilde\eta}]}{2\sqrt2}
\right)>0,
\]
where $\widetilde\eta$ is an independent copy of $\eta$.  Next set
\[
\theta:=\min\bigl(\delta^2/8,\delta\kappa_{\delta,\eta}\bigr),
\qquad
\lambda:=\min\bigl(\delta^3/2^{20},\delta\theta/2^{17}\bigr),
\qquad
s_0:=\frac{\theta^2}{2^{14}}.
\]
Let $S_\eta$ be the bounded linear probe from
Lemma~\ref{prop:exceptional-probes}. Choose a
constant $K_{\delta,\eta}$ so that
\[
K_{\delta,\eta}\ge \max\bigg\{\frac{4}{\delta^2},\norm{S_\eta}_{L^\infty}\bigg\} \ge 1,
\]
as $0<\delta<1$. In the case $\eta^2\not\equiv1$, let
$T_\eta$ be the bounded quadratic probe and increase $K_{\delta,\eta}$ (if necessary) to also satisfy
\[
K_{\delta,\eta}\ge \norm{T_\eta}_{L^\infty}.
\]
Define the head set $H$ by
\[
H(a,b):=\{1\}\cup\{i\ge2:\max(|a_i|,|b_i|)>\lambda\}.
\]
Then the cardinality of $H$ satisfies
\[
\#H\le 3+2\lambda^{-2} \eqqcolon d^{\ast}(\delta,\eta),
\]
and every tail coordinate $i\notin H$ satisfies
\[
\max(|a_i|,|b_i|)\le \lambda.
\]

The first step is to separate the proof into two mutually exclusive regimes.
Either one of the two signs makes the tail nearly one-dimensional, in which
case the argument is driven by head rigidity, or both signs retain substantial
tail mass, in which case anticoncentration becomes effective. We gather this simple case split as the following result, whose proof we omit. 

\begin{lemma}[Tail sign dichotomy]\label{lem:dichotomy}
Exactly one of the following alternatives holds:
\begin{enumerate}[label=\rm(\roman*)]
\item \textit{Diffuse:} We have
\[
\norm{(a-b)\mathbf 1_{H^c}}_2^2\ge16s_0
\quad\text{and}\quad
\norm{(a+b)\mathbf 1_{H^c}}_2^2\ge16s_0.
\]
\item \textit{Concentrated:} there exists $\tau\in\{\pm1\}$ such that
\[
\norm{(a-\tau b)\mathbf 1_{H^c}}_2<\frac{\theta}{32}.
\]
\end{enumerate}
\end{lemma}



\subsection{The concentrated branch}

We now turn to the first regime.  In the concentrated case, the tail of
$a-\tau b$ is small for a suitable sign $\tau$, so any substantial discrepancy
between the two vectors must be visible on the head.  The next proposition
shows that the quadratic observable $X^2-Y^2$ controls precisely those
head coefficients.

\begin{proposition}[Head entry bounds with one exceptional coordinate]
\label{prop:head-entry}
The following estimates hold:
\begin{enumerate}[label=\rm(\arabic*)]
\item For any indexes $i,j\in H\cap\{2,3,\dots\}$,
\[
|a_i a_j-b_i b_j|\le K_{\delta,\eta}\,\eps;
\]
\item For any index $j\in H\cap\{2,3,\dots\}$,
\[
|a_1a_j-b_1b_j|\le K_{\delta,\eta}\,\eps;
\]
\item if $\eta^2\not\equiv1$, then
\[
|a_1^2-b_1^2|\le K_{\delta,\eta}\,\eps.
\]
\end{enumerate}
\end{proposition}

\begin{proof}
For $i=j\ge2$, use the quadratic probe $T_i$ for $\xi_i$:
\[
|a_i^2-b_i^2|
=\bigl|\E[(X^2-Y^2)T_i]\bigr|
\le \norm{T_i}_{L^\infty}\eps
\le K_{\delta,\eta}\eps.
\]
For distinct $i,j\ge2$, use the linear probes $S_i,S_j$:
\[
2|a_i a_j-b_i b_j|
=\bigl|\E[(X^2-Y^2)S_iS_j]\bigr|
\le \norm{S_i}_{L^\infty}\norm{S_j}_{L^\infty}\eps
\le K_{\delta,\eta}\eps.
\]
This proves \rm(1).

For \rm(2), use the product $S_\eta S_j$.  Exactly as in the proof of
Lemma~\ref{lem:head-sign}, all terms vanish except the mixed
exceptional-good term, so
\[
2|a_1a_j-b_1b_j|
=\bigl|\E[(X^2-Y^2)S_\eta S_j]\bigr|
\le \norm{S_\eta}_{L^\infty}\norm{S_j}_{L^\infty}\eps
\le K_{\delta,\eta}\eps.
\]

Finally, if $\eta^2\not\equiv1$, use $T_\eta$:
\[
|a_1^2-b_1^2|
=\bigl|\E[(X^2-Y^2)T_\eta]\bigr|
\le \norm{T_\eta}_{L^\infty}\eps
\le K_{\delta,\eta}\eps.
\]
\end{proof}

To convert these entrywise bounds into a global contradiction, we use a basic fact about rank-one matrices. Intuitively,  it isolates the geometric fact that two orthogonal unit vectors cannot agree too well on the
head unless the corresponding head difference matrix is small.

\begin{lemma}[Rank-one inequality]\label{lem:rank-one}
For $x,y\in\R^m$, let $A(x,y)_{ij}=x_ix_j-y_iy_j$ be the difference matrix between $x$ and $y$. Then
\[
\norm{x-y}_2^2\,\norm{x+y}_2^2\le 2\norm{A(x,y)}_F^2.
\]
\end{lemma}

\begin{proof}
Writing $s=\norm{x}_2^2$, $t=\norm{y}_2^2$, and $c=\ip{x}{y}$, the left-hand
side equals $(s+t)^2-4c^2$, while the right-hand side is
$2(s^2+t^2-2c^2)$.  Thus the right-hand side has an extra factor $(s-t)^2$, which is nonnegative.
\end{proof}

Combining the head-entry bounds with the rank-one inequality yields the
concentrated branch.  The proof splits off the genuinely degenerate
Rademacher-only head, since only there can the exceptional coordinate fail to
provide a quadratic probe.

\begin{proposition}[Concentrated branch]\label{prop:concentrated}
If there exists $\tau\in\{\pm1\}$ such that
\[
\norm{(a-\tau b)\mathbf 1_{H^c}}_2\le \frac{\theta}{32},
\]
then there exists a constant $c_{\mathrm{line}}(\delta,\eta)>0$ such that
\[
\eps\ge c_{\mathrm{line}}(\delta,\eta).
\]
\end{proposition}

\begin{proof}
We distinguish two cases.

\smallskip
\noindent\emph{Case 1: $\eta$ is not Rademacher, or $H\cap\{2,3,\dots\}$ is
nonempty.}
In this case the head coefficients are rigidly linked by
Proposition~\ref{prop:head-entry}: either directly through the diagonal
estimate for the exceptional coordinate or, in the Rademacher case, through a
good head coordinate exactly as in the proof of Lemma~\ref{lem:head-sign}.  In
particular there exists a sign, which must be the same $\tau$, such that the
full head difference matrix of $a\mathbf 1_H$ and $\tau b\mathbf 1_H$ has all
entries bounded by $K_{\delta,\eta}\eps$. Thus, applying Lemma~\ref{lem:rank-one} for
\[
x:=a\mathbf 1_H,
\qquad
y:=\tau b\mathbf 1_H,
\]
we obtain 
\[
\norm{x-y}_2^2\,\norm{x+y}_2^2
\le 2\#H^2 K_{\delta,\eta}^2\eps^2.
\]
Since $(a-\tau b)\mathbf{1}_{H^c}$ has norm at most $\theta/32$ and
$\norm{a-\tau b}_2^2=2$ by orthogonality between $a$ and $b$, the head part satisfies
\[
\norm{x-y}_2^2
=2-\norm{(a-\tau b)\mathbf 1_{H^c}}_2^2
\ge 1.
\]
	Hence
	\[
	\norm{x+y}_2\le \sqrt{2}(\#H) K_{\delta,\eta}\eps
	\le \sqrt{2}d^*(\delta,\eta)K_{\delta,\eta}\eps.
	\]
	Set
	\[
	\alpha:=\norm{(a+\tau b)\mathbf 1_H}_2,
	\qquad
	\beta:=\norm{(a-\tau b)\mathbf 1_{H^c}}_2.
	\]
	Then
	\[
	\alpha\le \sqrt2\,d^*(\delta,\eta)K_{\delta,\eta}\eps,
	\qquad
	\beta\le \theta/32.
	\]
	We now spell out the gluing step.  Write
	\[
	p:=(a-\tau b)\mathbf 1_H,
	\qquad
	q:=(a-\tau b)\mathbf 1_{H^c},
	\]
	\[
	r:=(a+\tau b)\mathbf 1_H,
	\qquad
	s:=(a+\tau b)\mathbf 1_{H^c},
	\]
	and let
	\[
	P:=p_1\eta+\sum_{i\in H\cap\{2,3,\dots\}} p_i\xi_i,
	\qquad
	Q:=\sum_{i\notin H} q_i\xi_i,
	\]
	\[
	R:=r_1\eta+\sum_{i\in H\cap\{2,3,\dots\}} r_i\xi_i,
	\qquad
	S:=\sum_{i\notin H} s_i\xi_i.
	\]
	Then
	\[
	X-\tau Y=P+Q,
	\qquad
	X+\tau Y=R+S,
	\]
	so
	\[
	X^2-Y^2
	=(X-\tau Y)(X+\tau Y)
	=(P+Q)(R+S).
	\]
	Since the supports of $p,r$ and $q,s$ are disjoint, the pair $(P,R)$ is
	independent of the pair $(Q,S)$. In particular, $P$ and $S$ are independent.
	Therefore
	\[
	\abs{PS}\le \abs{X^2-Y^2}+\abs{PR}+\abs{QS}+\abs{QR},
	\]
	and hence
	\begin{equation}
    \label{eq:aux_conc_branch}
	\E[\abs{P}]\,\E[\abs{S}]
	=\E[\abs{PS}]
	\le
	\eps+\E[\abs{PR}]+\E[\abs{QS}]+\E[\abs{QR}].
	\end{equation}
	By Cauchy--Schwarz,
	\[
	\E[\abs{PR}]\le \norm{P}_{L^2}\norm{R}_{L^2},
	\qquad
	\E[\abs{QS}]\le \norm{Q}_{L^2}\norm{S}_{L^2},
	\qquad
	\E[\abs{QR}]\le \norm{Q}_{L^2}\norm{R}_{L^2}.
	\]
	Moreover,
	\[
	\norm{P}_{L^2}=\norm{p}_2=\sqrt{2-\beta^2},
	\qquad
	\norm{R}_{L^2}=\norm{r}_2=\alpha,
	\]
	\[
	\norm{Q}_{L^2}=\norm{q}_2=\beta,
	\qquad
	\norm{S}_{L^2}=\norm{s}_2=\sqrt{2-\alpha^2}.
	\]
	Assume for contradiction that $\eps<c_{\mathrm{line}}(\delta,\eta)$, where
	$c_{\mathrm{line}}(\delta,\eta)$ will be chosen so small that
	\[
	\sqrt2\,d^*(\delta,\eta)K_{\delta,\eta}c_{\mathrm{line}}(\delta,\eta)\le1.
	\]
	Then $\alpha\le1$ and
	\[
    1\le \min\{\|P\|_{L^2},\|S\|_{L^2}\}\le \max\{\|P\|_{L^2},\|S\|_{L^2}\} 
	\le \sqrt{2}
	\]
	Consequently from \eqref{eq:aux_conc_branch}
	\begin{equation}
    \label{eq:aux2_conc_branch}
	\E[\abs{P}]\,\E[\abs{S}]
	\le
	\eps+\sqrt2\,\alpha+\sqrt2\,\beta+\alpha\beta.
	\end{equation}
	
To lower-bound the left-hand side, we decompose $P$ as follows
	\[
	P=p_1\eta+G,
	\qquad \text{where} \quad 
	G:=\sum_{i\in H\cap\{2,3,\dots\}} p_i\xi_i.
	\]
	Either $\abs{p_1}\ge \norm{p}_2/\sqrt2$ or $\norm{G}_{L^2}\ge \norm{p}_2/\sqrt2$.
	In the first case, we apply Jensen's inequality
	\[
	\E[\abs{P}]
	=\E\bigl[\abs{p_1\eta+G}\bigr]
	\ge
	\abs{p_1}\,\E[\abs{\eta}]
	\ge
	\kappa_{\delta,\eta}\norm{p}_2.
	\]
	In the second case, let $G'$ be an independent copy of $G$.  Then
	\[
	2\E[\abs{P}]
	=\E\bigl[\abs{p_1\eta+G}+\abs{p_1\eta+G'}\bigr]
	\ge \E[\abs{G-G'}].
	\]
	Now
	\[
	G-G'
	=
	\sum_{i\in H\cap\{2,3,\dots\}} p_i(\xi_i-\xi_i'),
	\]
	where $(\xi_i')$ is an independent copy of $(\xi_i)$.  The variables
	$\xi_i-\xi_i'$ are independent and centered, satisfy
	\[
	\norm{\xi_i-\xi_i'}_{L^2}=\sqrt2,
	\qquad
	\norm{\xi_i-\xi_i'}_{L^1}\ge \delta,
	\]
	and therefore
	\[
	\norm{\xi_i-\xi_i'}_{L^1}
	\ge \frac{\delta}{\sqrt2}\norm{\xi_i-\xi_i'}_{L^2}.
	\]
	Applying Lemma~\ref{lem:l1lower},
	\[
	\E[\abs{G-G'}]
	\ge
	\frac{\delta/\sqrt2}{2\sqrt2}
	\Bigl(\sum_{i\in H\cap\{2,3,\dots\}} p_i^2
	\norm{\xi_i-\xi_i'}_{L^2}^2\Bigr)^{1/2}
	=\frac{\delta}{2}\norm{G}_{L^2}
	\ge 2\kappa_{\delta,\eta}\norm{p}_2.
	\]
	In both cases we obtain
	\[
	\E[\abs{P}]
	\ge \kappa_{\delta,\eta}\norm{p}_2
	\ge \kappa_{\delta,\eta}.
	\]
	Since $S$ is supported on $H^c$, only good coordinates enter it.  Thus
	Lemma~\ref{lem:l1lower} gives
	\[
	\E[\abs{S}]
	\ge \frac{\delta}{2\sqrt2}\norm{s}_2
	\ge \frac{\delta}{2\sqrt2}.
	\]
	Combining the last two displays with inequality \ref{eq:aux2_conc_branch}, we have
	\[
	\frac{\delta\kappa_{\delta,\eta}}{2\sqrt2}
	\le
	\eps+\sqrt2\,\alpha+\sqrt2\,\beta+\alpha\beta.
	\]
	Using the bounds for $\alpha$ and $\beta$, we infer that
	\[
	\frac{\delta\kappa_{\delta,\eta}}{2\sqrt2}
	\le
	\Bigl(1+2d^*(\delta,\eta)K_{\delta,\eta}
	+\frac{\theta}{16}d^*(\delta,\eta)K_{\delta,\eta}\Bigr)\eps
	+\frac{\sqrt2\,\theta}{32}.
	\]
	By the choice of $\theta$,
	\[
	\frac{\sqrt2\,\theta}{32}
	\le \frac{\delta\kappa_{\delta,\eta}}{8\sqrt2}.
	\]
	Hence, after decreasing $c_{\mathrm{line}}(\delta,\eta)$ if necessary so that
	\[
	\Bigl(1+2d^*(\delta,\eta)K_{\delta,\eta}
	+\frac{\theta}{16}d^*(\delta,\eta)K_{\delta,\eta}\Bigr)
	c_{\mathrm{line}}(\delta,\eta)
	\le \frac{\delta\kappa_{\delta,\eta}}{8\sqrt2},
	\]
	which is strictly smaller than
$\delta\kappa_{\delta,\eta}/(2\sqrt2)$ whenever
	$\eps<c_{\mathrm{line}}(\delta,\eta)$,thus reaching a contradiction. We conclude that
	$\eps\ge c_{\mathrm{line}}(\delta,\eta)$.

\smallskip
\noindent\emph{Case 2: $\eta$ is Rademacher and $H=\{1\}$.}
Write
\[
u:=a\mathbf 1_{H^c},
\qquad
v:=b\mathbf 1_{H^c},
\qquad
\rho:=\theta/32,
\]
and let
\[
U:=\sum_{i\ge2} a_i\xi_i,
\qquad
V:=\sum_{i\ge2} b_i\xi_i.
\]
Since $H=\{1\}$, every coordinate entering $U$ and $V$ is good.  The
concentrated assumption states precisely that
\[
\norm{u-\tau v}_2=\norm{U-\tau V}_{L^2}\le \rho.
\]

Set
\[
R:=U-\tau V,
\qquad
S:=U+\tau V,
\qquad
\alpha:=a_1-\tau b_1,
\qquad
\beta:=a_1+\tau b_1.
\]
Then
\[
X-\tau Y=\alpha\eta+R,
\qquad
X+\tau Y=\beta\eta+S.
\]
Because $\norm{a}_2=\norm{b}_2=1$ and $\ip{a}{b}=0$, we have
\[
\norm{a-\tau b}_2^2=\norm{a+\tau b}_2^2=2.
\]
Therefore
\[
\alpha^2+\norm{R}_{L^2}^2
=\alpha^2+\norm{u-\tau v}_2^2
=2,
\qquad
\beta^2+\norm{S}_{L^2}^2
=\beta^2+\norm{u+\tau v}_2^2
=2.
\]
In particular,
\[
|\alpha|\ge \sqrt{2-\rho^2}\ge1.
\]
Next, we estimate $\beta$. To this end, notice that
\[
|\beta|\le |\alpha\beta|
=|a_1^2-b_1^2|
=\bigl|\norm{v}_2^2-\norm{u}_2^2\bigr|
=\bigl|\langle u-\tau v,u+\tau v\rangle\bigr|,
\]
thus it follows that
\[
|\beta|
\le \norm{u-\tau v}_2\norm{u+\tau v}_2
\le \sqrt2\,\rho.
\]
Hence $\norm{S}_{L^2}^2 = \|u+\tau v\|_2^2\ge1$, and Lemma~\ref{lem:l1lower} gives
\[
\norm{S}_{L^1}
=\Bigl\|\sum_{i\ge2}(a_i+\tau b_i)\xi_i\Bigr\|_{L^1}
\ge \frac{\delta}{2\sqrt2}\norm{u+\tau v}_2
\ge \frac{\delta}{2\sqrt2}.
\]
On the other hand,
\[
\norm{R}_{L^1}\le \norm{R}_{L^2}\le \rho.
\]
Now condition on the sigma-algebra generated by the good coordinates.  Since
$\eta$ is an independent Rademacher variable,
\[
X^2-Y^2
=(X-\tau Y)(X+\tau Y)
=(\alpha\eta+R)(\beta\eta+S)
=\alpha\beta+RS+\eta(\alpha S+\beta R).
\]
Therefore
\[
\E\bigl[\abs{X^2-Y^2}\mid R,S\bigr]
=\frac12\Bigl(
\abs{\alpha\beta+RS+\alpha S+\beta R}
+
\abs{\alpha\beta+RS-\alpha S-\beta R}
\Bigr)
\]
\[
=\max\bigl(\abs{\alpha\beta+RS},\abs{\alpha S+\beta R}\bigr)
\ge \abs{\alpha S+\beta R}.
\]
Taking expectations on both sides, by the law of iterated expectation we have 
\begin{equation*}
\E[|X^2-Y^2|]\ge \|\alpha S+\beta R\|_{L^1}\ge  |\alpha|\|S\|_{L^1}-|\beta|\|R\|_{L^1}.
\end{equation*}
Recalling the definition of $\varepsilon$ and the estimates for the $L^1$ norm of $S$ and $R$, we conclude that
\[
\eps
\ge |\alpha|\norm{S}_{L^1}-|\beta|\norm{R}_{L^1}
\ge \frac{\delta}{2\sqrt2}-\sqrt2\,\rho^2.
\]
Since $\rho=\theta/32\le \delta^2/256$ and $0<\delta\le1$, the right-hand side
is bounded below by a positive constant depending only on $\delta$ (for
instance by $\delta/4$).  After decreasing
$c_{\mathrm{line}}(\delta,\eta)$ if necessary, this proves Case~2.
\end{proof}

\subsection{The diffuse branch}

We now move to the complementary regime.  Here both $a-b$ and $a+b$ retain a
definite amount of $\ell^2$ mass on the tail, and every tail coordinate is
individually small.  This is exactly the setting in which block
anticoncentration can be used.

Set
\[
c:=\frac{(a-b)\mathbf 1_{H^c}}{\sqrt2},
\qquad
d:=\frac{(a+b)\mathbf 1_{H^c}}{\sqrt2}.
\]
Because $1\in H$, the vectors $c$ and $d$ are supported only on good
coordinates.  In the diffuse case both have $\ell^2$ norm at least
$\sqrt{8s_0}$, while each coordinate is bounded by $\sqrt2\,\lambda$.

The next lemma provides the quantitative probabilistic input.  It says that
after arbitrary deterministic shifts, the two tail combinations cannot both
remain close to zero with high probability.
\begin{lemma}[Quantitative cross anticoncentration]\label{lem:cross}
There exists $r_*=r_*(\delta,\eta)>0$ such that, in the diffuse case,
\[
\Prob\bigl(\min(\abs{X_c+u},\abs{X_d+v})>r_*\bigr)\ge \frac14, \quad 
\qquad(u,v\in\R).
\]
\end{lemma}

\begin{proof}
We prove first a one-variable small-ball estimate.  Let $e\in \ell^2(\N)$ be
supported on $H^c$ that satisfies
\[
\norm{e}_2^2\ge 8s_0
\qquad \text{and} \quad
\sup_{i\ge2}\abs{e_i}\le \sqrt2\,\lambda.
\]
	We claim that there exists $r_0=r_0(\delta,\eta)>0$ such that
\begin{equation}\label{eq:onevar-smallball}
\Prob\bigl(\abs{X_e-t}\le r_0\bigr)\le \frac38
\qquad(t\in\R).
\end{equation}
Once this is proved, the lemma follows immediately by applying
\eqref{eq:onevar-smallball} to $e=c$ with $t=-u$ and to $e=d$ with $t=-v$.
Indeed, in the diffuse case both $c$ and $d$ satisfy the displayed
hypotheses, so
\[
\Prob\bigl(\abs{X_c+u}\le r_0\bigr)\le \frac38,
\qquad
\Prob\bigl(\abs{X_d+v}\le r_0\bigr)\le \frac38.
\]
By the union bound,
\[
\Prob\bigl(\abs{X_c+u}\le r_0\ \text{or}\ \abs{X_d+v}\le r_0\bigr)
\le \frac34,
\]
and therefore
\[
\Prob\bigl(\min(\abs{X_c+u},\abs{X_d+v})>r_0\bigr)\ge \frac14.
\]
Thus it remains to establish the claim \eqref{eq:onevar-smallball}.

\smallskip
\noindent\emph{Step 1: a block partition.}
Set
\[
M:=\Bigl\lceil 2^{16}\delta^{-2}\Bigr\rceil,
\qquad
\eta:=\frac{s_0}{2M},
\qquad
r_0:=\frac{\delta\sqrt{\eta}}{16\sqrt{2}}.
\]
	Because $\delta\le1$ implies
	\[
	M\le 2^{17}\delta^{-2},
	\]
	and because
	\[
	2\lambda^2
	\le \frac{\delta^2\theta^2}{2^{33}}
	\le \frac{\theta^2}{2^{15}M}
	=\eta,
	\]
	every coordinate satisfies $e_i^2\le \eta$.  Since
\[
\norm{e}_2^2\ge 8s_0>2M\eta=s_0,
\]
we may choose pairwise disjoint finite sets $G_1,\dots,G_M\subset H^c$ such
that
\begin{equation}\label{eq:block-energy}
\eta\le \sum_{i\in G_j} e_i^2\le 2\eta
\qquad(1\le j\le M).
\end{equation}
Indeed, construct the blocks greedily.  Suppose $G_1,\dots,G_{j-1}$ have
already been chosen.  Their total energy is at most $2(j-1)\eta$, so the
unused part of $e$ still has energy at least
\[
8s_0-2(j-1)\eta\ge 8s_0-2M\eta=7s_0> \eta.
\]
Hence one can choose a finite subset of the remaining support whose energy is
at least $\eta$, and by minimality its energy is at most $\eta+\sup_i e_i^2
\le 2\eta$.  This produces the next block.

\smallskip
\noindent\emph{Step 2: symmetrized block sums have uniform two-sided mass.}
Let $X_e'$ be an independent copy of $X_e$, realized on a second copy of the
probability space.  For each $i\ge2$ define
\[
D_i:=\xi_i(\omega_1)-\xi_i'(\omega_2).
\]
Then the variables $(D_i)_{i\ge2}$ are independent, centered, and symmetric,
with
\[
\norm{D_i}_{L^2}^2
=\E[(\xi_i-\xi_i')^2]
=2.
\]
Moreover, by Jensen's inequality applied conditionally on $\xi_i$,
\[
\E[\abs{D_i}]
=\E\bigl[\E[\abs{\xi_i-\xi_i'}\mid \xi_i]\bigr]
\ge \E\bigl[\abs{\xi_i-\E[\xi_i']} \bigr]
=\E[\abs{\xi_i}]
\ge \delta.
\]
Hence
\[
\E[\abs{D_i}]
\ge \frac{\delta}{\sqrt2}\norm{D_i}_{L^2}.
\]
For each block define
\[
Y_j:=\sum_{i\in G_j} e_iD_i,
\qquad
\sigma_j^2:=\sum_{i\in G_j} e_i^2.
\]
By \eqref{eq:block-energy},
\[
\eta\le \sigma_j^2\le 2\eta.
\]
Applying Lemma~\ref{lem:l1lower} to the independent centered family
$(D_i)_{i\in G_j}$, with lower $L^1$ ratio $\delta/\sqrt2$, gives
\[
\E[\abs{Y_j}]
\ge
\frac{(\delta/\sqrt2)}{2\sqrt2}
\Bigl(\sum_{i\in G_j} e_i^2\norm{D_i}_{L^2}^2\Bigr)^{1/2}
=\frac{\delta}{4}\bigl(2\sigma_j^2\bigr)^{1/2}
=\frac{\delta}{2\sqrt{2}}\sigma_j
\ge \frac{\delta\sqrt{\eta}}{2\sqrt{2}}.
\]
Set
\[
a_0:=\frac{\delta\sqrt{\eta}}{4\sqrt{2}}=4r_0,
\]
and notice that $a_0\le \E[\abs{Y_j}]/2$.  We define the event 
\[
A_j:=\{\abs{Y_j}\ge a_0\}.
\]
Since
\[
\E[Y_j^2]=2\sigma_j^2\le 4\eta,
\]
we obtain
\[
\E[\abs{Y_j}]
\le a_0+\E[\abs{Y_j}\mathbf 1_{A_j}]
\le a_0+\Prob(A_j)^{1/2}\E[Y_j^2]^{1/2}\le a_0 +\Prob(A_j)^{1/2}(4\eta)^{1/2} .
\]
Because $\E[\abs{Y_j}]\ge 2a_0$, the second term on the right-hand side satisfies
\[
\Prob(A_j)^{1/2}\,(4\eta)^{1/2}\ge a_0.
\]
Therefore
\begin{equation}\label{eq:block-active}
\Prob(A_j)\ge \frac{a_0^2}{4\eta}=\frac{\delta^2}{128}.
\end{equation}

\noindent\emph{Step 3: many active blocks with high probability.}
Let
\[
N:=\sum_{j=1}^M \mathbf 1_{A_j},
\]
and notice that it is a sum of independent random variables because the blocks are
disjoint.  By \eqref{eq:block-active},
\[
\E[N]
=\sum_{j=1}^M \Prob(A_j)
\ge M\frac{\delta^2}{128}
\ge 512.
\]
Since $\operatorname{Var}(\mathbf 1_{A_j})\le \E[\mathbf 1_{A_j}]$, the independence implies that
\[
\operatorname{Var}(N) = \sum_{j=1}^M \operatorname{Var}(\mathbf 1_{A_j})\le \sum_{j=1}^M \E[\mathbf 1_{A_j}]=  \E[N].
\]
Chebyshev's inequality gives
\[
\Prob(N<256)
\le
\Prob\bigl(\abs{N-\E[N]}\ge \E[N]-256\bigr)
\le
\frac{\operatorname{Var}(N)}{(\E[N]-256)^2}
\le
\frac{\E[N]}{(\E[N]-256)^2}
\le \frac{512}{256^2}
<\frac1{16}.
\]

\smallskip
\noindent\emph{Step 4: a Sperner bound for the symmetrized sum.}
Let
\[
U:=\sum_{j=1}^M Y_j.
\]
We claim that for every $s\in\R$,
\begin{equation}\label{eq:symm-smallball}
\Prob(\abs{U-s}\le 2r_0)\le \frac18.
\end{equation}
Indeed,
\[
\Prob(\abs{U-s}\le 2r_0)
\le
\Prob(\abs{U-s}\le 2r_0,\ N\ge256)+\Prob(N<256).
\]
We have already bounded the second term by $1/16$, so it remains to estimate
the first one.

Condition on the absolute values $(\abs{Y_j})_{j=1}^M$.  Since each $Y_j$ is
symmetric, conditional on $(\abs{Y_j})_j$ the signs of the nonzero $Y_j$ are
independent Rademacher variables.  On the event $\{N\ge256\}$ let
\[
J:=\{j:\abs{Y_j}\ge a_0\},
\qquad
m:=\#J\ge256.
\]
After freezing the inactive coordinates and the magnitudes $\abs{Y_j}$, the
condition $\abs{U-s}\le 2r_0$ becomes
\[
\sum_{j\in J}\varepsilon_j w_j \in I,
\]
where $w_j:=\abs{Y_j}\ge a_0$ and $I$ is an interval of diameter
$4r_0=a_0$.  Since $a_0<2w_j$ for every $j\in J$, the family of sign patterns
for which the weighted sum falls in $I$ is an antichain in $\{\pm1\}^J$.
Indeed, if two such sign patterns were comparable in the natural subset order,
their sums would differ by at least $2w_j\ge2a_0>a_0$, which is impossible
inside an interval of diameter $a_0$.  By Sperner's theorem, the conditional
probability of this event is therefore at most
\[
\frac{\binom{m}{\lfloor m/2\rfloor}}{2^m}
\le \frac1{\sqrt m}
\le \frac1{16},
\]
because $m\ge256$.  Averaging over the conditioning yields
\[
\Prob(\abs{U-s}\le 2r_0,\ N\ge256)\le \frac1{16}.
\]
Together with the bound for $\Prob(N<256)$, this proves
\eqref{eq:symm-smallball}.

\smallskip
\noindent\emph{Step 5: unsymmetrization.}
Let
\[
R:=\sum_{i\notin G_1\cup\cdots\cup G_M} e_iD_i,
\]
and notice that $R$ is independent of $U$, because it is built from coordinates disjoint
from those entering the block sums.  Also by construction, we have
\[
X_e(\omega_1)-X_e'(\omega_2)=U+R.
\]
Fix $t\in\R$.  If both $\abs{X_e(\omega_1)-t}\le r_0$ and
$\abs{X_e'(\omega_2)-t}\le r_0$, then
\[
\abs{X_e(\omega_1)-X_e'(\omega_2)}\le 2r_0.
\]
Hence since $X_{e}'$ is an independent copy of $X_e$,
\[
\Prob(\abs{X_e-t}\le r_0)^2 = \Prob(\abs{X_e-t}\le r_0\cap\abs{X_e'-t}\le r_0)
\le
\Prob(\abs{X_e(\omega_1)-X_e(\omega_2)}\le 2r_0)
=
\Prob(\abs{U+R}\le 2r_0).
\]
Conditioning on $R$ and using \eqref{eq:symm-smallball},
\[
\Prob(\abs{U+R}\le 2r_0)
\le
\sup_{s\in\R}\Prob(\abs{U-s}\le 2r_0)
\le \frac18.
\]
Therefore
\[
\Prob(\abs{X_e-t}\le r_0)\le \frac1{\sqrt8}<\frac38.
\]
This proves \eqref{eq:onevar-smallball}, and with it the lemma.
\end{proof}

Once this shifted anticoncentration estimate is available, the lower bound for
$\eps$ follows from the same factorization of $X^2-Y^2$ that was
already used in the compactness proof.

\begin{proposition}[Diffuse branch]\label{prop:diffuse}
If the diffuse alternative of Lemma~\ref{lem:dichotomy} holds, then
\[
\eps\ge c_{\mathrm{cross}}(\delta,\eta)>0.
\]
\end{proposition}

\begin{proof}
Let
\[
h_\pm:=X_{(a\pm b)\mathbf 1_H}.
\]
Then
\[
X^2-Y^2
=2\bigl(X_d+h_+/\sqrt2\bigr)\bigl(X_c+h_-/\sqrt2\bigr),
\]
and therefore
\[
\abs{X^2-Y^2}
\ge
2\min\Bigl(\abs{X_d+h_+/\sqrt2},\abs{X_c+h_-/\sqrt2}\Bigr)^2.
\]
Condition on the head variables $h_\pm$ and apply Lemma~\ref{lem:cross} with
the shifts $u=h_-/\sqrt2$ and $v=h_+/\sqrt2$.  This yields a uniform positive
lower bound for the conditional expectation of the squared minimum, hence for
$\eps$.  The resulting constant depends only on $\delta$ and $\eta$.
\end{proof}

\subsection{Conclusion of the second proof}

At this point the proof is immediate: the dichotomy shows that every
orthogonal pair falls into exactly one of the two branches, and each branch
gives a quantitative lower bound for the same quantity $\eps$.

\begin{proof}[Second proof of Theorem~\ref{thm:main}]
By Proposition~\ref{prop:orth-reduction}, it is enough to prove a uniform
orthogonal lower bound.  Let $a,b$ be orthogonal unit vectors and form $X$
and $Y$ as above.  Apply the dichotomy of Lemma~\ref{lem:dichotomy}:
if the concentrated alternative holds, then
Proposition~\ref{prop:concentrated} yields
\[
\eps\ge c_{\mathrm{line}}(\delta,\eta).
\]
If the diffuse alternative holds, then Proposition~\ref{prop:diffuse} yields
\[
\eps\ge c_{\mathrm{cross}}(\delta,\eta).
\]
Since $\eps\le 2\Delta$, we conclude that
\[
\Delta
=\norm{|X|-|Y|}_{L^2}
\ge
\frac12\min\bigl(c_{\mathrm{line}}(\delta,\eta),c_{\mathrm{cross}}(\delta,\eta)\bigr)
=:c_{\delta,\eta}>0.
\]
The orthogonal reduction then gives
\[
\min_{\sigma\in\{\pm1\}}\norm{X-\sigma Y}_{L^2}
\le \frac{\sqrt2}{c_{\delta,\eta}}\norm{|X|-|Y|}_{L^2}
\]
for all $X,Y\in \overline{\Span}(\eta,\xi_i:i\ge2)$, completing the proof of
Theorem~\ref{thm:main}.
\end{proof}

\section{Comments and Open Problems}\label{sec:comments}

The compactness and quantitative proofs suggest several natural directions for
further work.  Although the theorem proved here settles the qualitative
one-exception problem in the real Hilbertian setting, the mechanisms that enter
the proof are flexible enough that one is led almost immediately to adjacent
questions. 


\begin{question}
Classical phase retrieval concerns recovery from the absolute value modulo the group symmetry $\{\pm1\}$.  It is
natural to ask: are there analogous theorems for other nonlinearities modulo possibly larger compact symmetry groups?  %{\color{red} Maybe mention declipping?}
\end{question}

The compactness
proof is particularly well suited to the above direction, since it separates the
problem into an exact limit identity and a geometric classification of the
corresponding support set.  It points toward a broader geometric framework in which the
support theorem, rather than the specific real-sign symmetry, is the central
object.

The next question concerns the extent to which the present theorem can be
sharpened, and how much should persist in a wider Banach-function-space
setting.

\begin{question}
Which Banach-function-space analogues of the present theorem remain true once
the ambient space is no longer Hilbertian? In particular, is there an $L^p$-analogue of our result when $p<2$? The answer is yes when $p\geq 2$,
as an immediate consequence of the interpolation/extrapolation theory developed in \cite{FOTP}.%The orthogonal reduction is
%special to Hilbert spaces, so a different mechanism would be needed.
\end{question}

This problem seems especially interesting, as the function-space formulation
of stable phase retrieval is not inherently Hilbertian, even though the present
proof is.  Any successful extension would likely require replacing
exact orthogonality by a more flexible geometric notion while retaining enough
structure to control the dichotomy between near and far vectors.
\begin{question}
To what extent can independence be weakened in the above theorems?
%special to Hilbert spaces, so a different mechanism would be needed.
\end{question}
The motivation for the above problem is that there are many examples of sequences of functions (for example, sparse Fourier series) that behave in certain senses like independent random variables. To what extent can the methods of this paper be generalized to cover these examples as well?

\section{Appendix: Discussions on the LEAN verification}

Theorem XXX, which is the main technical result of this manuscript has been Lean-verified. The proof has been mostly generated using LLMs (GPT 5.4 through codex and Claude Opus 4.6 through Claude Code, using the lean-lsp-mcp and custom skills). 

This section discusses the formalization of the statement, which can be found in the file {\scriptsize \texttt{showcase.lean}}. This file was written using LLMs and then heavily modified by the authors for simplicity. At the definition level, this file depends only on Mathlib. In particular, reviewing the statements in this file is enough to attest the semantic correctness of the Lean-verified statement, conditional on the semantics of the Mathlib library itself. Moreover, the statements in this file are simple and we hope that, at a broad level, can be understood by mathematicians without Lean knowledge. This section attempts to explain the contents of the file.

We start by defining $\Omega$ as a probability space with probability measure $\mathbb P$. In Mathlib, this is expressed as
\begin{leancode}
variable { Ω : Type*} [MeasureSpace Ω] [IsProbabilityMeasure ( ℙ : Measure Ω)]
\end{leancode}

One can then define stability constant by
\begin{leancode}
def StabilityConstant (A B : ℝ) : ℝ :=
  let dStar := 2 + 2 ^ 41 / A ^ 6
  let K := max (4 / A ^ 2) (2 / (1 - B))
  let cLine := A ^ 2 / (2 ^ 8  * dStar * K)
  let cCross  := A ^ 8  / 2  ^ 57 
  (1 / 2 : ℝ) * min cCross cLine
\end{leancode}

This is enough to define the main theorem
\begin{leancode}
theorem quantitative_orthogonal_lower_bound
    {n : ℕ}
    (ξ : Fin n → Ω → ℝ)--                        Let ξ₁, …, ξₙ : Ω → ℝ be random variables
    (hMeas : ∀ i, Measurable (ξ i))--            (measurable)
    (hIndep : iIndepFun--                        that are jointly independent.
      (m := fun _ => inferInstance) ξ ℙ)
    (hMean : ∀ i, !$\mathbb E$![ξ i] = 0)--                  Assume that !$\mathbb E$![ξᵢ] = 0
    (hVar : ∀ i, (!$\mathbb E$![(ξ i) ^ 2] : ℝ) = 1)--       and !$\mathbb E$![ξᵢ²] = 1.
    {A B : ℝ} (hA : 0 < A) (hB : B < 1)--        Let 0 < A, B < 1
    (hL1_lower : ∀ i, A ≤ !$\mathbb E$![|ξ i|])--            such that A ≤ !$\mathbb E$![|ξᵢ|],
    (hL1_upper : ∀ i, i.val ≠ 0 → !$\mathbb E$![|ξ i|] ≤ B)--and !$\mathbb E$![|ξᵢ|] ≤ B for i ≠ 0.
    (a b : Fin n → ℝ)--                          Let (a₁,... aₙ), (b₁,... aₙ) be real vectors
    (haNorm : ∑ i, a i ^ 2 = 1)--                satisfying ‖a‖₂ = 1,
    (hbNorm : ∑ i, b i ^ 2 = 1)--                ‖b‖₂ = 1,
    (hab : ∑ i, a i * b i = 0)--                 and ⟨a, b⟩ = 0.
    : --                                         Then,
    StabilityConstant A B ≤ --                   StabilityConstant(A,B) ≤ √ !$\mathbb E$![(|Xa| − |Xb|)²]
      Real.sqrt (!$\mathbb E$![(|∑ i, a i • ξ i| - |∑ i, b i • ξ i|) ^ 2])
\end{leancode}

\begin{enumerate}
    \item {\scriptsize \texttt{Fin n}} represents a finite type of $n$ elements, which can be mapped into the real numbers using the ``.val" map. This is Lean's representation of $[n]=\{0,\dots, n-1\}$. A finite family of random variables $\xi_1, \dots \xi_n$ is represented as a map from {\scriptsize \texttt{Fin n}} to the type of random variables.
    \item The type of random variables itself, is the set of functions $\Omega \to \mathbb R$, with the hypotheses that they are measurable. In particular, a list of random variables has type $[n]\to(\Omega \to \mathbb R$))
    \item Independence of (families of) random variables is defined at the level of \textit{functions} mapping from an index set to  random variables. This is defined in Mathlib via {\scriptsize \texttt{indepFun}}, which takes three inputs: 
    \begin{enumerate}
        \item The $\sigma$-algebras to be considered in the \textit{target} space of the random variables (in this case, automatically inferred to be the Borel sigma algebra over the reals by {\scriptsize \texttt{inferInstance}}.
        \item The family of random variables (in this case $\xi$)
        \item The probability measure (in this case $\mathbb P$)
    \end{enumerate}
    \item Hypotheses and \textit{data} that is used to state the theorem is written in either parentheses or brackets. This has no difference from the perspective of stating the result, and is only a simplification when later using the theorem: Types in brackets can be \textit{automatically inferred} from the context by Lean, while types in parentheses will have to be explicitly provided by the user.
    
\end{enumerate}

\begin{remark}[On Lean and division]
Division has two particularities in Lean that warrant a remark.

\begin{enumerate}
    \item Division between natural numbers is Euclidean. That is, as an
    identity in $\mathbb N$, Lean interprets 
    \[
        2 / 8 = 0.
    \]
    By contrast, division over the reals is the usual real division. To make sure that
    Lean reads a natural number as a real number, we include the symbol
    {\scriptsize\texttt{(2 : $\mathbb R$)}}. Thus, the division above becomes
    \[
        (2 / 8 : \mathbb R) = 1/4.
    \]
    This kind of type cast appears in the last line of the
    definition of {\scriptsize\texttt{StabilityConstant}}.

    \item Even for real numbers, division is a total operation in Lean: every
    expression of the form $x/y$ has a value, including when $y=0$. The
    Mathlib convention is
    \[
        x / 0 = 0.
    \]
    This does not mean that division by zero is mathematically harmless.
    Rather, it means that Lean will allow the expression syntactically, and
    the responsibility is on the proof to establish that the relevant
    denominator is nonzero whenever the intended mathematical argument requires
    it. In our case, our theorem is meaningful as long as the denominators in {\scriptsize\texttt{StabilityConstant}} are nonzero (which holds because $0<A$ and $0< B < 1$).
\end{enumerate}
\end{remark}

\bibliographystyle{plain}
\bibliography{refs}

\end{document}


\nocite{calderbank2022stable}
\nocite{bandeira2014saving}
\nocite{eldar2014stability}
\nocite{cahill2016infinite}
\nocite{AG-continuous-frames}
\nocite{alaifari2021gabor}
\nocite{grohs2020survey}
\nocite{FOTP}
\nocite{christ2022examples}
\nocite{localphase}
\nocite{chen2019phase}
\nocite{MR3746047}
\nocite{KS}
\nocite{DB}
\nocite{MR3260258}
\nocite{MR3069958}
\nocite{MR4298599}
\nocite{sanz1984phase}
\nocite{balan2013}
\nocite{bartusel2021phase}
\nocite{casazza2017weak}
\nocite{casazza2017norm}
\nocite{casazza2013}
\nocite{AG-continuous-frames}
\nocite{fg22}
\nocite{MR3901674}
\nocite{grohs2021stable}
\nocite{luef2021gaussian}
\nocite{kashin2021sampling}
\nocite{LT21}
\nocite{DPSTT}
\nocite{NOU}
\nocite{akrami2021note}
\nocite{MR3367650}
\nocite{MR3275625}
\nocite{CCS-JMAA}
\nocite{CDOS}
\nocite{CHL}
\nocite{FG-atomic}
%\nocite{MR2011364}
\nocite{MR0423039}
\nocite{FOTP}